% ============================================
% The Opus Audio Codec: Theory and Implementation
% Educational manual for RFC 6716 and the ruopus pure-Rust implementation
% ============================================

\section{Introduction}\label{sec:introduction}

This manual provides a comprehensive overview of the algorithms,
optimizations, and design decisions that make up the Opus audio
codec~\cite{rfc6716}. The focus is on understanding both the theoretical
foundations and the practical implementations of the techniques Opus combines:
the information theory of its entropy coder, the linear prediction of its speech
path, the transform coding of its music path, and the fixed-point and SIMD
machinery that lets all of it run in real time.

The material emphasizes the \termdef{why} and \termdef{how} of each component:
\emph{why} certain design choices matter, and \emph{how} they are implemented
efficiently. It covers both the mathematics and the practical computational
considerations that make these methods work in production systems.

Opus is the IETF's general-purpose, royalty-free audio codec, standardized in
2012 and updated by RFC 8251~\cite{rfc8251}. It is the default audio codec of
WebRTC and is widely deployed in streaming and conferencing. Its defining
characteristic is that it is not one codec but two, combined behind a single
bitstream. A linear-predictive speech coder (\termdef{SILK}) and a
transform-domain music coder (\termdef{CELT}, the Constrained-Energy Lapped
Transform) share one entropy coder, one packet format, and one container. The
encoder selects, per frame, between SILK, CELT, or both together (the
\termdef{hybrid} mode). The combination spans a range no single technique covers
well: 6~kb/s narrowband speech to 510~kb/s fullband stereo music, frame sizes
from 2.5~ms to 60~ms, and sample rates from 8~kHz to 48~kHz, all decoded by one
receiver that is not told in advance which mode each packet uses.

\subsection{The Reference Implementation}\label{subsec:reference_impl}

Throughout, the implementation under discussion is
\texttt{ruopus}, a pure-Rust realization of RFC 6716 with zero \emph{required}
runtime dependencies (the default build pulls in one optional FFT crate for
faster decoding; \code{default-features = false} is fully dependency-free).
Every code listing in this manual is drawn from that source. Three of its design
choices shape the discussion:

\begin{itemize}[itemsep=3pt]
  \item \textbf{Bit-exact by construction.} The bitstream is defined by the
    range coder (\cref{sec:rangecoder}), not by any one encoder's choices. The
    implementation's entropy arithmetic follows the RFC byte-for-byte, and the
    decoder is validated against the official test vectors on the
    \emph{final-range oracle}, a per-packet checksum of the entropy-coder
    state that must agree with the reference to the bit.
  \item \textbf{Safe by default.} The crate denies \code{unsafe} everywhere
    except a handful of explicitly annotated SIMD kernels
    (\cref{sec:optimizations}). The signal-processing math is ordinary safe
    Rust; the only machine intrinsics sit in encoder hot loops where a
    faster, approximate result is provably harmless (\cref{sec:optimizations}).
  \item \textbf{Decode-first.} The decoder is the normative half of the
    specification and the half the conformance vectors exercise. The encoder is
    free to make any choice the decoder can follow; this asymmetry recurs
    throughout the manual and is the key to several optimizations.
\end{itemize}

\subsection{Target Audience and Prerequisites}\label{subsec:prerequisites}

This material assumes familiarity with:
\begin{itemize}[itemsep=3pt]
  \item Discrete signals, sampling, and the discrete Fourier transform
  \item Linear algebra (vectors, matrices, inner products, projection)
  \item Basic probability and the idea of entropy / information content
  \item Programming and algorithm analysis; the listings are in Rust, but no
    Rust expertise is required to follow them
\end{itemize}

No prior knowledge of speech coding, transform coding, arithmetic coding, or
psychoacoustics is assumed; each is introduced where it is first needed. A
reader who has worked through a companion treatment of FFT-based spectral
analysis will find the transform-coding sections (CELT) immediately familiar,
and the linear-prediction sections (SILK) a complementary story told in the
time domain.

\subsection{Scope and Organization}\label{subsec:organization}

The manual is organized bottom-up, from the bit to the file, mirroring the
layering of the codec itself:

\begin{enumerate}[itemsep=3pt]
  \item \textbf{Architecture Overview} (\cref{sec:architecture}) --- the two
    codecs, three modes, the shared range coder, and the sample-rate model.
  \item \textbf{The Range Coder} (\cref{sec:rangecoder}) --- entropy coding,
    arithmetic/range coding, the ICDF representation, raw bits, \code{tell()},
    and the final-range oracle.
  \item \textbf{Packet Framing} (\cref{sec:packet}) --- the TOC byte, the 32
    configurations, frame packing codes 0--3, padding, and validity rules.
  \item \textbf{CELT: Transform Coding} (\cref{sec:celt}) --- the MDCT, the
    Vorbis window and TDAC, pyramid vector quantization, CWRS enumeration, band
    splitting, energy coding, bit allocation, and the encoder's analysis.
  \item \textbf{SILK: Linear-Predictive Coding} (\cref{sec:silk}) --- linear
    prediction, LSF/NLSF quantization, long-term (pitch) prediction, the
    noise-shaping quantizer, gains, the shell coder, and resampling.
  \item \textbf{The Opus Layer} (\cref{sec:opus_layer}) --- mode dispatch,
    the SILK/CELT crossover, redundancy at transitions, packet-loss
    concealment, in-band FEC, and discontinuous transmission.
  \item \textbf{Ogg Encapsulation} (\cref{sec:ogg}) --- pages, the Ogg CRC,
    lacing, packet reassembly, \code{OpusHead}/\code{OpusTags}, and granule
    timing.
  \item \textbf{Multistream and Surround} (\cref{sec:multistream}) --- coupled
    and uncoupled streams, channel mapping families.
  \item \textbf{Computational Optimizations} (\cref{sec:optimizations}) ---
    SIMD dot products, the SIMD PVQ search, the MDCT pre-rotation, the FFT
    backend seam, table-driven CWRS, and buffer reuse.
  \item \textbf{Numerical Considerations and Conformance}
    (\cref{sec:numerical}) --- fixed-point versus float, bit-exactness, the
    oracle, and the conformance criterion.
\end{enumerate}

\subsection{Notation Conventions}\label{subsec:notation}

Throughout this manual:
\begin{itemize}[itemsep=2pt]
  \item $x[n]$ denotes a time-domain sample at index $n$; $X[k]$ a
    frequency-domain coefficient at bin $k$
  \item $f_s$ denotes a sample rate in Hz; Opus's internal reference rate is
    always 48~kHz
  \item $N$ denotes a transform length or vector dimension; $K$ a pulse count
  \item $\mathrm{LM}$ (``log multiplier'') denotes $\log_2$ of the frame-size
    multiplier: $\mathrm{LM}=0,1,2,3$ for 2.5, 5, 10, 20~ms CELT frames
  \item The range coder's state is the pair $(\mathit{val}, \mathit{rng})$
  \item \termdef{\Q{n}} denotes a fixed-point value scaled by $2^n$: a \Q{15}
    value $v$ represents the real number $v / 2^{15}$
  \item $\sgn(x)$ is the sign of $x$; $\floor*{\cdot}$ and $\ceil*{\cdot}$ are
    floor and ceiling
  \item Section references like ``\S4.3.4'' without further qualification refer
    to RFC 6716~\cite{rfc6716}
\end{itemize}

\section{Architecture Overview}\label{sec:architecture}

This section describes the codec as a whole: its components, the path a packet
takes through them, and the design decisions that the rest of the manual
elaborates. Later sections treat each layer in detail.

\subsection{Two Codecs, One Bitstream}\label{subsec:two_codecs}

Opus is built on the observation that audio divides into two regimes that reward
opposite coding strategies. The codec implements both and selects between them
per frame.

\begin{concept}{Speech versus Music: Two Coding Strategies}
  \textbf{Speech} is produced by a physical system (the vocal folds exciting
  the vocal tract) that is well modeled as a slowly varying \emph{linear
  filter} driven by a sparse \emph{excitation}. Coding speech efficiently means
  coding that model: estimate the filter, estimate the pitch, and transmit only
  the residual that the model fails to predict. This is
  \termdef{linear-predictive
  coding} (LPC), and it is the approach SILK takes, in the \emph{time
  domain}, at
  internal rates of 8, 12, or 16~kHz.

  \textbf{Music} is arbitrary: dense, polyphonic, and percussive, with energy
  spread across the full 20~kHz audible band and no simple source model. Coding
  music efficiently means working in the \emph{frequency domain}, where the
  ear's own frequency analysis allows bits to be spent where they are audible
  and withheld where masking hides them. This is \termdef{transform coding}, and
  it is the approach CELT takes, via the Modified Discrete Cosine Transform
  (MDCT) at 48~kHz.

  Neither technique performs both tasks well: LPC degrades on polyphonic
  material, and transform coding spends bits inefficiently on the strong
  predictable structure of speech. Opus therefore retains both and switches
  between them per frame; in the overlap region (wideband speech with a musical
  background) it runs them together as the \termdef{hybrid} mode.
\end{concept}

This yields the three \termdef{modes}, selected by the packet header
(\cref{sec:packet}) and dispatched by the decoder (\cref{sec:opus_layer}):

\begin{center}
  \begin{tabular}{lll}
    \toprule
    \textbf{Mode} & \textbf{Coder(s)} & \textbf{Typical use} \\
    \midrule
    SILK-only & SILK & low-bitrate speech, 8--16~kHz audio band \\
    Hybrid    & SILK (low) + CELT (high) & wideband speech + full bandwidth \\
    CELT-only & CELT & music, low delay, high bitrate \\
    \bottomrule
  \end{tabular}
\end{center}

In hybrid mode the split is by frequency: SILK codes the signal up to roughly
8~kHz (it runs at a 16~kHz internal rate), and CELT codes the band above,
starting at MDCT band 17 (\cref{subsec:hybrid}). Both coders write into the same
range coder, so there is no boundary marker between them in the byte stream: the
CELT data is simply a continuation of the same range-coded message that began
with the SILK data.

\subsection{The Layered Stack}\label{subsec:stack}

\Cref{fig:stack-table} lists the layers as the implementation modularizes them,
bottom to top. The lower layers (range coder, packet, Ogg) are pure framing and
entropy: exact integer arithmetic, \code{no\_std}-compatible, with no audio
math. The middle layers (SILK, CELT) perform the signal processing. The top
layer (the Opus decoder and encoder) provides the glue: mode dispatch,
resampling, and the cross-frame state used to handle mode transitions.

\begin{table}[h]
  \centering
  \begin{tabular}{lll}
    \toprule
    \textbf{Layer} & \textbf{Specification} & \textbf{Responsibility} \\
    \midrule
    \code{ogg}         & RFC 3533 / 7845 & file container, pages, timing \\
    \code{multistream} & RFC 7845 \S5.1.1 & surround: $N$
    coupled/mono streams \\
    \code{decoder}/\code{encoder} & RFC 6716 \S4/\S5 & mode dispatch,
    transitions, PLC/FEC/DTX \\
    \code{silk}        & RFC 6716 \S4.2 & linear-predictive speech coder \\
    \code{celt}        & RFC 6716 \S4.3 & MDCT transform coder \\
    \code{packet}      & RFC 6716 \S3 & TOC byte, frame packing \\
    \code{range}       & RFC 6716 \S4.1/\S5.1 & entropy coder \\
    \bottomrule
  \end{tabular}
  \caption{The Opus layer stack, bottom to top.}
  \label{fig:stack-table}
\end{table}

\subsection{Decoding a Packet: An Overview}\label{subsec:packet_life}

To decode one Opus packet, the decoder performs the following steps:

\begin{enumerate}[itemsep=3pt]
  \item Reads the first byte, the \termdef{TOC} (Table of Contents), which
    names the mode, audio bandwidth, frame size, and channel count
    (\cref{sec:packet}).
  \item Splits the remaining bytes into 1--48 \emph{frames} according to the
    TOC's frame-packing code.
  \item For each frame, constructs a range decoder over its bytes and, depending
    on the mode, runs the SILK decoder, the CELT decoder, or both into a shared
    range decoder.
  \item Resamples the decoder's internal output (always referenced to 48~kHz) to
    the caller's requested rate, and applies any cross-frame smoothing for mode
    transitions and redundancy (\cref{sec:opus_layer}).
\end{enumerate}

The entry point parses the packet and hands it to a per-frame driver, itself a
long chain of mode- and bandwidth-conditioned branches rather than one tidy
dispatch, since transitions and hybrid redundancy each need their own
handling:

\begin{lstlisting}[language=Rust]
pub fn decode_packet(&mut self, data: &[u8]) -> Result<Vec<f32>, PacketError> {
    // Empty or TOC-only packet: conceal one frame (DTX / loss).
    if data.len() <= 1 {
        return Ok(self.decode_lost(self.last_frame_size / self.ds));
    }
    let packet = Packet::parse(data)?;
    Ok(self.decode_parsed(&packet))
}

fn decode_parsed(&mut self, packet: &Packet) -> Vec<f32> {
    let toc = packet.toc();
    self.stream_channels = usize::from(toc.channels());
    let mut out = Vec::new();
    for frame in packet.frames() {
        let pcm = self.decode_frame(
            frame, toc.mode(), toc.bandwidth(),
            toc.frame_size().samples_per_channel_48k(),
        );
        out.extend_from_slice(&pcm);
    }
    out
}
\end{lstlisting}

\subsection{Sample Rates and the 48 kHz Reference}\label{subsec:internal_rate}

The relationship between sample rates in Opus is a common source of confusion.
The model is stated once here and referred to throughout.

\begin{concept}{The 48 kHz Reference Clock}
  Opus has three distinct notions of sample rate:
  \begin{itemize}[itemsep=2pt]
    \item The \textbf{audio bandwidth} (narrowband, mediumband, wideband,
      super-wideband, fullband) is a \emph{property of the content}: how much of
      the spectrum carries signal. It is named in the TOC and caps the highest
      coded frequency (4, 6, 8, 12, or 20~kHz).
    \item The \textbf{internal rate} is the rate a coder actually runs at. SILK
      runs at 8, 12, or 16~kHz; CELT always runs at 48~kHz.
    \item The \textbf{API rate} is the rate the caller asks for (8, 12, 16,
      24, or 48~kHz), at which PCM is delivered.
  \end{itemize}

  Crucially, \textbf{all timing is expressed in 48~kHz samples} regardless of
  the other two. A 20~ms frame is always 960 samples ``at 48~kHz''; a
  2.5~ms frame is 120. Granule positions in the Ogg container
  (\cref{sec:ogg}) count 48~kHz samples. The decoder resamples from the internal
  rate up to 48~kHz (CELT is already there), then \emph{down} to the API rate by
  an integer factor $\mathit{ds} = 48000 / f_s$. This single reference clock is
  what lets SILK frames, CELT frames, and the container agree on ``when'' a
  sample occurs even though they compute it at different rates.
\end{concept}

\subsection{From Quantities to Symbols}\label{subsec:bits_overview}

Every numeric quantity Opus transmits (a gain, an LSF index, a pulse vector,
an energy delta) becomes a symbol in the range coder. The encoder and decoder
share static probability models (\termdef{ICDF} tables, \cref{sec:rangecoder})
for the structured fields and fall back to \termdef{raw bits} for the
near-uniform ones. This arrangement has an important consequence:

\begin{importantnote}
  The bitstream is defined by the \emph{sequence of range-coder symbols}, not
  by the floating-point intermediate values the encoder computed to choose them.
  Two encoders that pick the same symbols produce identical bytes, even if one
  used SIMD and \code{rsqrt} approximations and the other used scalar
  double-precision. This is why the decoder must be bit-exact (it is normative)
  but the encoder need only be \emph{valid} (any choice the decoder can follow
  is legal). Most of \cref{sec:optimizations} exploits exactly this freedom.
\end{importantnote}

\section{The Range Coder}\label{sec:rangecoder}

Every bit Opus transmits passes through the range coder. Because the bitstream
is defined entirely by the sequence of range-coder symbols, and because
bit-exact conformance is verified through the coder's internal state
(\cref{subsec:oracle}), it is the natural starting point and is treated here
first and in detail.

\subsection{Why Entropy Coding}\label{subsec:why_entropy}

Suppose a symbol $s$ occurs with probability $p(s)$. Information theory says the
\emph{ideal} cost of transmitting it is its \termdef{self-information},
$-\log_2 p(s)$ bits~\cite{shannon1948mathematical,cover2006elements}. A common
symbol ($p$ near 1) should cost a fraction of a bit; a rare one, many bits. Over
a long message drawn from a known distribution, the average cost approaches the
\termdef{entropy} $H = -\sum_s p(s)\log_2 p(s)$, and no lossless code can do
better.

A fixed-length code cannot reach this bound: it must spend a whole number of
bits per symbol, wasting the fractional part. \termdef{Arithmetic coding}, and
its integer-friendly cousin \termdef{range coding}~\cite{martin1979range,
witten1987arithmetic}, do reach it, by abandoning the idea that each symbol maps
to its own bit pattern. Instead, the \emph{entire message} is encoded as a
single number in $[0, 1)$, and each symbol narrows the interval that number must
lie in.

\begin{concept}{Range Coding by Interval Subdivision}
  Begin with the interval $[0, 1)$. To encode a symbol $s$ whose cumulative
  distribution carves out the sub-range $[\mathit{fl}/\mathit{ft},\,
  \mathit{fh}/\mathit{ft})$ (low and high cumulative frequencies out of a total
  $\mathit{ft}$), \emph{replace} the current interval with that fraction of
  itself:
  \begin{equation}\label{eq:rc_subdivide}
    [\,\ell, \ell + r\,) \;\longmapsto\;
    \left[\,\ell + r\tfrac{\mathit{fl}}{\mathit{ft}},\;
    \ell + r\tfrac{\mathit{fh}}{\mathit{ft}}\,\right).
  \end{equation}
  The interval shrinks by a factor of $(\mathit{fh}-\mathit{fl})/\mathit{ft} =
  p(s)$ at each step, so after the whole message its width is $\prod_i p(s_i) =
  2^{-\sum_i -\log_2 p(s_i)}$. Transmitting \emph{any} number inside the final
  interval takes about $-\log_2(\text{width}) = \sum_i -\log_2 p(s_i)$ bits,
  the entropy. The decoder, knowing the same model, replays the subdivisions:
  it sees which sub-range the received number falls in, emits that symbol,
  rescales, and repeats.
\end{concept}

The practical obstacle is that the interval bounds need ever more precision as
they shrink, which is impossible with finite-width integers. The fix, due to
Pasco and refined by Martin, is \termdef{renormalization}: as soon as the
interval is narrow enough that its top bits are settled, \emph{emit those bits}
and shift them out, restoring headroom. Opus renormalizes a byte at a time.

\subsection{The Integer Range Coder}\label{subsec:integer_rc}

Opus's range coder (\S4.1 for the decoder, \S5.1 for the encoder) works in
32-bit integers. Its state is two words:

\begin{itemize}[itemsep=2pt]
  \item \code{rng}: the current interval width, kept in
    $(2^{23}, 2^{31}]$.
  \item \code{val}: on the encoder, the interval's low end; on the decoder,
    the offset of the received number from the top of the interval.
\end{itemize}

The governing constants name the renormalization threshold and the byte width:

\begin{lstlisting}[language=Rust]
const SYM_BITS:  u32 = 8;            // emit/consume one byte at a time
const SYM_MAX:   u32 = 255;          // 2^8 - 1
const CODE_BITS: u32 = 32;           // width of the (val, rng) state
const CODE_TOP:  u32 = 1 << 31;      // 2^31: initial encoder range
const CODE_BOT:  u32 = 1 << 23;      // 2^23: renormalize when rng falls to here
const CODE_SHIFT: u32 = 23;          // CODE_BITS - SYM_BITS - 1; carry bit position

#[inline]
const fn ilog(x: u32) -> u32 { u32::BITS - x.leading_zeros() }  // 1 + floor(log2 x)
\end{lstlisting}

The lower bound $\mathit{rng} > 2^{23}$ guarantees that the division
$\mathit{rng}/\mathit{ft}$ in \cref{eq:rc_subdivide} keeps enough significant
bits for the largest total frequency Opus uses ($\mathit{ft} \le 2^{16}$). Note
\code{ilog} is just the hardware count-leading-zeros instruction; it appears
everywhere the coder needs $\floor*{\log_2}$.

\subsubsection{Decoding a symbol}

Decoding is two steps: \code{decode} reports \emph{where} the received value
sits in $[0, \mathit{ft})$, the caller maps that to a symbol, then \code{update}
shrinks the state to that symbol's sub-range and renormalizes.

\begin{lstlisting}[language=Rust]
pub fn decode(&mut self, ft: u32) -> u32 {
    // Locate val within [0, ft): the symbol s satisfies fl <= decode() < fh.
    ft - (self.val / (self.rng / ft) + 1).min(ft)
}

pub fn update(&mut self, fl: u32, fh: u32, ft: u32) {
    let s = self.rng / ft;
    self.val -= s * (ft - fh);            // shift past the symbols above this one
    self.rng = if fl > 0 { s * (fh - fl) } // interior symbol: exact sub-width
               else       { self.rng - s * (ft - fh) }; // first symbol keeps remainder
    self.normalize();
}
\end{lstlisting}

Two details are easy to miss and important to get right. First, the first symbol
($\mathit{fl}=0$) is given the \emph{remainder} of the range
($\mathit{rng} - s\,(\mathit{ft}-\mathit{fh})$), not
$s\cdot\mathit{fh}$; this is
how the coder accounts for the truncation in $s = \floor*{\mathit{rng}/
\mathit{ft}}$ and stays reversible. Second, \code{decode} caps its result at
$\mathit{ft}$ via \code{.min(ft)} to defend against a corrupt stream pushing the
quotient out of range.

\subsubsection{The renormalization loop}

When \code{rng} drops to $2^{23}$ or below, the decoder shifts in a fresh byte:

\begin{lstlisting}[language=Rust]
fn normalize(&mut self) {
    while self.rng <= CODE_BOT {
        self.nbits_total += SYM_BITS;
        self.rng <<= SYM_BITS;
        // Reassemble the next 8-bit symbol from the leftover bit of the
        // previous byte (high) and the top 7 bits of the new byte (low).
        let byte = self.read_byte();
        let sym = (self.leftover_bit << 7) | (byte >> 1);
        self.leftover_bit = byte & 1;
        self.val = ((self.val << SYM_BITS) + (SYM_MAX - sym)) & 0x7FFF_FFFF;
    }
}
\end{lstlisting}

The one-bit shimmy (\code{leftover\_bit}) is a peculiarity of Opus's coder: the
decoder works with a 7-bits-from-this-byte, 1-bit-carried representation so that
the encoder's carry propagation (below) lands on byte boundaries. The
\code{\& 0x7FFF\_FFFF} keeps \code{val} within 31 bits.

\subsubsection{Encoding a symbol}

The encoder mirrors the decoder exactly; it must, or the two will disagree
and the stream will be undecodable:

\begin{lstlisting}[language=Rust]
pub fn encode(&mut self, fl: u32, fh: u32, ft: u32) {
    let r = self.rng / ft;
    if fl > 0 {
        self.val += self.rng - r * (ft - fl);
        self.rng  = r * (fh - fl);
    } else {
        self.rng -= r * (ft - fh);
    }
    self.normalize();
}
\end{lstlisting}

The encoder's \code{normalize} emits the top byte of \code{val} instead of
reading one, but is otherwise the same loop, with one complication that
deserves its own treatment.

\subsection{Carry Propagation}\label{subsec:carry}

When the encoder adds to \code{val} (the $\mathit{fl}>0$ branch above), the
addition can overflow into bits it has \emph{already emitted}. A byte written
out as \code{0xFF} might need to become \code{0x00} with a carry rippling into
the byte before it. The encoder cannot un-write bytes, so it does not write them
eagerly. Instead it buffers:

\begin{lstlisting}[language=Rust]
fn carry_out(&mut self, c: u32) {
    if c == SYM_MAX {
        self.ext += 1;                 // a run of 0xFF: defer, might all carry
    } else {
        let carry = (c >> SYM_BITS) as u8;
        if let Some(rem) = self.rem.take() {
            self.write_byte(rem + carry);     // resolve the buffered byte
        }
        if self.ext > 0 {
            // Flush the deferred 0xFF run, +carry turns them into 0x00.
            let sym = (SYM_MAX + u32::from(carry)) as u8;
            for _ in 0..self.ext { self.write_byte(sym); }
            self.ext = 0;
        }
        self.rem = Some((c & SYM_MAX) as u8);
    }
}
\end{lstlisting}

\begin{concept}{Why a 0xFF Run Must Be Held}
  A carry into a \code{0xFF} byte turns it into \code{0x00} and \emph{propagates
  further left}. A carry into any other byte simply increments it and stops. So
  the encoder can safely emit a non-\code{0xFF} byte only once it
  knows whether a
  carry is coming, which it learns from the \emph{next} byte. It therefore
  keeps one pending byte in \code{rem} and counts a run of \code{0xFF}s in
  \code{ext}. When a non-\code{0xFF} byte arrives with carry $c$: the pending
  byte becomes $\mathit{rem}+c$, and the held \code{0xFF}s become either
  \code{0xFF} (no carry) or \code{0x00} (carry). This buffering is invisible to
  the decoder; it simply reads the bytes that finally emerge. The mechanism is
  what makes the encoder/decoder \emph{symmetric to the bit}, which is the
  property the conformance oracle (\cref{subsec:oracle}) checks.
\end{concept}

\subsection{The ICDF Representation}\label{subsec:icdf}

Structured fields are coded against static probability tables. Opus stores them
not as a PDF or a forward CDF but as an \termdef{inverse cumulative distribution
function} (ICDF): for a model with total $2^{\mathit{ftb}}$, the table holds
$\mathit{icdf}[k] = 2^{\mathit{ftb}} - \mathit{fh}[k]$, the cumulative
\emph{tail} above each symbol, decreasing to a terminating zero.

\begin{concept}{Why Inverse CDFs}
  Two reasons. First, the decode primitive in \cref{eq:rc_subdivide} naturally
  produces ``how far down from the top is \code{val}'', which is exactly an
  upper-tail comparison, so an ICDF lets the decoder walk the table with one
  multiply and one compare per step and no running subtraction. Second, the
  tables are tiny (a handful of bytes), and storing the tail lets the common
  ``most probable symbol first'' ordering terminate early. The forward
  arithmetic ($\mathit{fl}, \mathit{fh}$) is recovered implicitly from adjacent
  ICDF entries.
\end{concept}

\begin{lstlisting}[language=Rust]
pub fn decode_icdf(&mut self, icdf: &[u8], ftb: u32) -> usize {
    let d = self.val;
    let r = self.rng >> ftb;          // ft is a power of two: shift, don't divide
    let mut k = 0usize;
    let mut t = self.rng;
    let mut s = r * u32::from(icdf[0]);
    while d < s {                     // walk down the tail until val sits in [s, t)
        k += 1;
        t = s;
        s = r * u32::from(icdf[k]);
    }
    self.val = d - s;                 // fold update() into the same step
    self.rng = t - s;
    self.normalize();
    k
}
\end{lstlisting}

Because $\mathit{ft} = 2^{\mathit{ftb}}$ is a power of two, the division
$\mathit{rng}/\mathit{ft}$ becomes a shift, a small but pervasive saving,
since ICDF decoding is the hottest path in SILK and in CELT's side information.

\subsection{Binary Symbols and Raw Bits}\label{subsec:binary_raw}

Two specializations round out the coder.

\textbf{Binary symbols with a power-of-two probability.} Many flags have a
probability of the form $2^{-\mathit{logp}}$. These skip the table entirely:

\begin{lstlisting}[language=Rust]
pub fn decode_bit_logp(&mut self, logp: u32) -> bool {
    let s = self.rng >> logp;        // sub-range for the "1" outcome
    let bit = self.val < s;
    if !bit { self.val -= s; }
    self.rng = if bit { s } else { self.rng - s };
    self.normalize();
    bit
}
\end{lstlisting}

\textbf{Raw bits.} Some quantities are effectively uniform (a sign, the low
bits of a pulse split), and modeling them as probabilities wastes effort.
Opus codes these as \termdef{raw bits}, written at the \emph{opposite end} of
the same buffer from the arithmetic-coded data.

\begin{concept}{The Two-Ended Buffer}
  A range-coded packet has \emph{two} write heads in one byte array. The
  arithmetic coder fills from the front, most-significant first, byte by byte as
  it renormalizes. Raw bits fill from the \emph{back}, least-significant first.
  The two heads grow toward each other; the packet is well-formed as long as
  they do not cross. This costs nothing: raw bits need no probability model,
  so they need not interleave with the arithmetic stream, and it lets the
  finalizer (\cref{subsec:tell}) pack the last arithmetic byte and the last raw
  byte into the same physical byte when they happen to fit. The decoder reads
  raw bits backwards from the end of the packet through a small 32-bit window,
  refilling a byte at a time.
\end{concept}

\begin{lstlisting}[language=Rust]
pub fn decode_raw_bits(&mut self, bits: u32) -> u32 {
    if self.nend_bits < bits {                       // refill the window from the end
        loop {
            self.end_window |= self.read_byte_from_end() << self.nend_bits;
            self.nend_bits += SYM_BITS;
            if self.nend_bits > WINDOW_SIZE - SYM_BITS { break; }
        }
    }
    let ret = self.end_window & ((1 << bits) - 1);    // LSB-first
    self.end_window >>= bits;
    self.nend_bits  -= bits;
    self.nbits_total += bits;                         // tell()/tell_frac() count these too
    ret
}
\end{lstlisting}

\textbf{Uniform integers} combine the two: a value in $[0, \mathit{ft})$ with
$\mathit{ft}$ not a power of two is split into an arithmetic-coded high part
(8 bits of model) and a raw-bit low part.

\begin{lstlisting}[language=Rust]
pub fn encode_uint(&mut self, t: u32, ft: u32) {
    let ftb = ilog(ft - 1);
    if ftb <= 8 {
        self.encode(t, t + 1, ft);                    // small: pure arithmetic
    } else {
        let ft_hi = ((ft - 1) >> (ftb - 8)) + 1;
        self.encode(t >> (ftb - 8), (t >> (ftb - 8)) + 1, ft_hi);
        self.encode_raw_bits(t & ((1 << (ftb - 8)) - 1), ftb - 8);  // low bits raw
    }
}
\end{lstlisting}

\subsection{Bit Budgeting: tell() and tell\_frac()}\label{subsec:tell}

Opus's bit allocation (\cref{subsec:celt_allocation}) needs to know, mid-frame,
exactly how many bits have been spent so far, and it needs that figure to
\emph{agree} between encoder and decoder, since both run the same allocation
arithmetic from the same running total. The coder exposes this through
\code{tell} and its eighth-of-a-bit refinement \code{tell\_frac}.

\begin{lstlisting}[language=Rust]
pub fn tell(&self) -> u32 { self.nbits_total - ilog(self.rng) }

pub(crate) fn tell_frac(nbits_total: u32, rng: u32) -> u32 {
    let mut lg = ilog(rng);
    let mut r = rng >> (lg - 16);     // Q15 fractional part of log2(rng)
    for _ in 0..3 {                   // 3 iterations -> 1/8-bit resolution
        r = (r * r) >> 15;
        let bit = r >> 16;
        lg = 2 * lg + bit;
        r >>= bit;
    }
    nbits_total * 8 - lg
}
\end{lstlisting}

\begin{concept}{Bits Used: Total Shifted In Minus Range Headroom}
  The intuition: \code{nbits\_total} counts every bit shifted into the coder.
  But the coder has not \emph{committed} all of them: the current range
  \code{rng} still spans $\log_2(\mathit{rng})$ bits of undecided value. So the
  bits actually consumed are \code{nbits\_total} $- \log_2(\mathit{rng})$.
  \code{tell()} uses the integer $\floor*{\log_2}$ (a conservative whole-bit
  count); \code{tell\_frac()} refines $\log_2(\mathit{rng})$ to 3
  fractional bits
  by the Newton-style squaring loop above, giving the $1/8$-bit
  resolution Opus's
  allocator works in. The two readouts agree on encoder and decoder because they
  are computed from the same $(\mathit{nbits\_total}, \mathit{rng})$.
\end{concept}

\subsection{The Final-Range Oracle}\label{subsec:oracle}

The range coder provides a single, powerful correctness check for any Opus
implementation, described below.

\begin{concept}{rng-Agreement and the Final Range}
  After encoding (or decoding) the same sequence of symbols, the encoder and
  decoder must hold \emph{identical} \code{rng} values, and, transitively,
  identical \code{val} and identical output bytes. RFC 6716 elevates this to a
  conformance tool: \code{OPUS\_GET\_FINAL\_RANGE} returns the
  coder's \code{rng}
  at the end of each packet, and a correct decoder's final range must equal the
  reference encoder's, \emph{exactly}, for every packet.

  This is a strong check. The entire decode (thousands of symbols across SILK
  and CELT, every table lookup, every renormalization) collapses into a single
  32-bit number that is wrong if \emph{any} step diverged. The
  \texttt{ruopus} conformance suite enforces it on all twelve official
  vectors: a passing decode is not merely close to the reference, it is
  arithmetically identical through the entropy coder. The PCM may differ by a
  rounding bit in the final float MDCT; the \emph{range} may not differ at all.
\end{concept}

In hybrid and redundancy cases (\cref{sec:opus_layer}) the packet carries more
than one range-coded segment; the implementation XORs their final ranges
together to form the single reported value, matching the reference's
bookkeeping.

\section{Packet Framing}\label{sec:packet}

The range coder produces the bytes of a single frame. The packet layer
(\S3) is the envelope around those bytes: it names how to interpret the frame,
allows several frames to share one packet, and defines the rules a decoder uses
to reject malformed input. It contains no audio math and no entropy coding, only
parsing.

\subsection{The Self-Delimiting Question}\label{subsec:self_delimit}

A design decision shapes the entire packet format: an Opus packet is \emph{not}
self-delimiting. It does not record its own length. Opus assumes a transport
layer (Ogg, RTP, WebRTC, Matroska) that already knows where each packet ends,
and declines to duplicate that information.

\begin{concept}{Why Packets Carry No Length Field}
  Storing a length in every packet would waste two or three bytes on data the
  transport already has. At the low bitrates Opus targets, a 20~ms packet can be
  under 20 bytes, so a redundant length prefix is a 10--15\% overhead. Opus
  instead infers frame boundaries from the total packet size supplied by the
  container. The cost is that the codec cannot be handed a bare byte stream and
  asked to find packet boundaries; it must be told, per packet, how many bytes
  belong to it. The one exception is the \termdef{self-delimiting} framing used
  inside multistream packets (\cref{sec:multistream}), where several
  Opus packets
  are concatenated and all but the last must announce their own length.
\end{concept}

\subsection{The TOC Byte}\label{subsec:toc}

Every packet begins with one \termdef{Table of Contents} (TOC) byte. It packs
three fields that, together, determine everything the decoder needs to dispatch
the packet:

\begin{equation}\label{eq:toc_layout}
  \underbrace{\texttt{c c c c c}}_{\text{config (5)}}\;
  \underbrace{\texttt{s}}_{\text{stereo (1)}}\;
  \underbrace{\texttt{n n}}_{\text{count code (2)}}
\end{equation}

The top five bits select one of 32 \termdef{configurations}; the next bit is the
mono/stereo flag; the low two bits are the frame-packing code
(\cref{subsec:framecodes}). The implementation treats the TOC as a transparent
wrapper over the byte, with each field a small accessor:

\begin{lstlisting}[language=Rust]
pub struct Toc(u8);

impl Toc {
    pub const fn config(self) -> u8 { self.0 >> 3 }            // bits 7..3
    pub const fn stereo(self) -> bool { (self.0 >> 2) & 1 == 1 } // bit 2
    pub const fn channels(self) -> u8 { 1 + ((self.0 >> 2) & 1) }
    pub const fn frame_count_code(self) -> u8 { self.0 & 0x3 } // bits 1..0
    pub const fn mode(self) -> Mode { /* from config */ }
    pub const fn bandwidth(self) -> Bandwidth { /* from config */ }
    pub const fn frame_size(self) -> FrameSize { /* from config */ }
}
\end{lstlisting}

\subsubsection{The 32 configurations}

The configuration number is not an arbitrary index; it is a structured encoding
of (mode, bandwidth, frame size). The 32 values partition into the three modes,
and within each mode the bandwidths and frame sizes are laid out in contiguous
runs (\cref{tab:configs}). This regularity lets the encoder compute a config as
a base value plus a frame-size offset, which is exactly how the encoder builds
its TOC (\cref{sec:opus_layer}).

\begin{table}[h]
  \centering
  \begin{tabular}{llll}
    \toprule
    \textbf{Config} & \textbf{Mode} & \textbf{Bandwidth} &
    \textbf{Frame sizes (ms)} \\
    \midrule
    0--3   & SILK-only & narrowband (NB)      & 10, 20, 40, 60 \\
    4--7   & SILK-only & mediumband (MB)      & 10, 20, 40, 60 \\
    8--11  & SILK-only & wideband (WB)        & 10, 20, 40, 60 \\
    12--13 & Hybrid    & super-wideband (SWB) & 10, 20 \\
    14--15 & Hybrid    & fullband (FB)        & 10, 20 \\
    16--19 & CELT-only & narrowband (NB)      & 2.5, 5, 10, 20 \\
    20--23 & CELT-only & wideband (WB)        & 2.5, 5, 10, 20 \\
    24--27 & CELT-only & super-wideband (SWB) & 2.5, 5, 10, 20 \\
    28--31 & CELT-only & fullband (FB)        & 2.5, 5, 10, 20 \\
    \bottomrule
  \end{tabular}
  \caption{The 32 TOC configurations (\S3.1, Table 2). The mode dictates which
    coder runs; the bandwidth caps the highest coded frequency; the
    frame size is
  always expressed at the 48~kHz reference.}
  \label{tab:configs}
\end{table}

Three properties of \cref{tab:configs} are worth noting. SILK-only modes support
the long 40 and 60~ms frames that suit speech but never the short 2.5 and 5~ms
frames; CELT-only modes are the reverse, supporting short low-delay frames down
to 2.5~ms. Hybrid mode exists only at the two highest bandwidths, because it is
defined as ``SILK wideband plus a CELT extension''; below super-wideband there
is no high band for CELT to add. And mediumband appears only for SILK, as CELT's
band layout has no 6~kHz cutoff.

\subsection{Frame-Packing Codes}\label{subsec:framecodes}

The two low bits of the TOC select one of four ways to pack frames into the
remaining bytes. The motivation is amortization: a single packet can carry up to
120~ms of audio (for example three 40~ms frames), spreading the one-byte TOC
overhead across multiple frames and reducing per-frame container overhead.

\begin{concept}{The Four Frame-Packing Codes}
  Let the bytes after the TOC be the \emph{payload}. The count code determines
  how the payload splits into frames:
  \begin{itemize}[itemsep=3pt]
    \item \textbf{Code 0:} one frame, occupying the entire payload.
    \item \textbf{Code 1:} two frames of \emph{equal} size; the payload length
      must be even, and each frame is half.
    \item \textbf{Code 2:} two frames of possibly \emph{unequal} size; an
      explicit length precedes the first frame, and the second takes the rest.
    \item \textbf{Code 3:} an arbitrary count $M$ of frames (1--48), with a
      following byte that selects constant-bitrate (CBR, all frames equal) or
      variable-bitrate (VBR, explicit lengths) packing, and optional padding.
  \end{itemize}
\end{concept}

Frame lengths, where coded explicitly (code 2 and VBR code 3), use a one- or
two-byte representation: a byte below 252 is the length itself; a byte of 252 to
255 is combined with a following byte as $\text{second} \times 4 +
\text{first}$, reaching up to 1275, the maximum frame size. A coded length of
zero is legal and denotes a dropped or DTX (discontinuous transmission) frame
that the decoder conceals.

The parser is a direct transcription of these four cases:

\begin{lstlisting}[language=Rust]
match toc.frame_count_code() {
    0 => { check_frame_len(rest.len())?; frames.push(rest); }
    1 => {
        if rest.len() % 2 != 0 { return Err(PacketError::Code1UnevenPayload); }
        let (a, b) = rest.split_at(rest.len() / 2);
        frames.push(a); frames.push(b);
    }
    2 => {
        let n1 = read_frame_len(&mut rest)?;     // explicit first length
        let frame1 = take(&mut rest, n1)?;
        frames.push(frame1); frames.push(rest);  // second takes the remainder
    }
    _ => { /* code 3: count byte, optional padding, CBR or VBR lengths */ }
}
\end{lstlisting}

\subsubsection{Code 3 in detail}

Code 3 carries a second header byte after the TOC: bit 7 is the VBR flag, bit 6
signals the presence of padding, and the low six bits give the frame count $M$.
Padding, when present, is a length encoded as a run of bytes (a value of 255
means ``254 bytes of padding, continue reading''), allowing a packet to be
inflated to a target size for constant-rate transport without affecting the
audio. In VBR packing the header is followed by $M-1$ explicit frame lengths,
with the final frame taking the remainder; in CBR packing all $M$ frames share
the payload equally, and a payload not divisible by $M$ is an error.

\subsection{Validity Rules}\label{subsec:validity}

RFC 6716 \S3.4 enumerates seven conditions, labeled [R1] through [R7], that a
well-formed packet must satisfy. They exist because a decoder fed hostile or
corrupt input must reject it cleanly rather than read out of bounds or allocate
unboundedly. The implementation's error type documents which rule(s) each
variant enforces (\cref{tab:rules}); several rules share an error because they
are checked by the same code path, and R4/R6/R7 each contribute to more than
one variant.

\begin{table}[h]
  \centering
  \begin{tabularx}{\textwidth}{lcX}
    \toprule
    \textbf{Error} & \textbf{Rule(s)} & \textbf{Condition} \\
    \midrule
    \code{Empty} & R1 & packet is at least one byte (has a TOC) \\
    \code{FrameTooLarge} & R2 & every frame is at most 1275 bytes \\
    \code{Code1UnevenPayload} & R3 & code 1 payload length is even \\
    \code{InvalidFrameLength} & R4, R7 & a coded frame length is truncated or
    overruns the payload \\
    \code{InvalidFrameCount} & R5 & code 3 count $M\ge 1$ and total $\le$
    120~ms \\
    \code{InvalidPadding} & R6, R7 & padding overruns the payload \\
    \code{CbrPayloadNotDivisible} & R6 & a CBR (code 2/3, non-VBR) payload
    does not divide evenly by the frame count \\
    \bottomrule
  \end{tabularx}
  \caption{The malformed-packet rules [R1]--[R7] (\S3.4) and the typed errors
  that enforce them.}
  \label{tab:rules}
\end{table}

\begin{importantnote}
  Rules R5 and R6 are the codec's defense against denial-of-service input. R5
  caps a packet at 120~ms of audio, so a single packet cannot demand an
  unbounded decode; the implementation checks $M \times \text{(tenths of ms per
  frame)} \le 1200$. R6 ensures the padding count and the CBR frame division,
  both attacker-controlled, cannot point past the end of the buffer or hide a
  malformed split. The packet rules are the first line that stops such input
  from reaching the signal-processing layers at all; the decoder is also
  exercised by differential fuzzing against the reference implementation
  (\cref{subsec:conformance}) to catch the arithmetic edge cases the rules
  alone do not.
\end{importantnote}

\section{CELT: Transform Coding}\label{sec:celt}

CELT, the Constrained-Energy Lapped Transform (\S4.3), is the music half of
Opus. It is a transform codec: it maps blocks of audio into the frequency domain
with a Modified Discrete Cosine Transform, then spends bits on the resulting
coefficients according to a psychoacoustic budget. Its distinguishing feature,
the one its name advertises, is that it codes the \emph{energy} of
each frequency
band explicitly and separately from the band's \emph{shape}, which keeps the
reconstructed spectrum's energy envelope faithful even when bits are scarce.

This section follows the decoder's data flow: the transform
(\cref{subsec:mdct}), the energy envelope (\cref{subsec:celt_energy}), the
pyramid vector quantizer that codes band shapes (\cref{subsec:pvq}), the
combinatorial enumeration that turns a quantized shape into an index
(\cref{subsec:cwrs}), the bit allocation that ties them together
(\cref{subsec:celt_allocation}), the band-coding machinery
(\cref{subsec:celt_bands}), and finally the encoder's analysis
(\cref{subsec:celt_encoder}).

\subsection{The MDCT and Lapped Transforms}\label{subsec:mdct}

A transform codec cannot simply cut the signal into blocks and transform each
independently: the discontinuities at block boundaries would produce audible
clicks, and the rectangular truncation would smear energy across frequencies
(spectral leakage). The standard remedy is to overlap consecutive blocks and
window them, but naive overlap doubles the data rate. The Modified Discrete
Cosine Transform~\cite{princen1986analysis,princen1987subband} resolves this
tension: it overlaps blocks by 50\% yet remains \emph{critically sampled},
producing exactly $N$ coefficients for every $N$ new input samples.

\begin{concept}{The Modified Discrete Cosine Transform}
  The MDCT maps $2N$ real input samples to $N$ real coefficients:
  \begin{equation}\label{eq:mdct}
    X[k] = \sum_{n=0}^{2N-1} x[n]\,
    \cos\!\left[\frac{\pi}{N}\left(n + \tfrac{1}{2} + \tfrac{N}{2}\right)
    \left(k + \tfrac{1}{2}\right)\right],
    \quad k = 0,\dots,N-1.
  \end{equation}
  The inverse (IMDCT) produces $2N$ samples from the $N$ coefficients. Because
  the forward transform discards half its inputs' worth of degrees of freedom,
  a single IMDCT cannot recover its input; the missing information is
  supplied by
  the \emph{overlapping} neighbor block. Each output block overlaps
  its successor
  by $N$ samples, and the two contributions are added. This is
  \termdef{time-domain aliasing cancellation} (TDAC): the IMDCT of one block
  contains a time-reversed alias of the input, and the overlapping block's alias
  is equal and opposite, so the sum cancels it exactly.
\end{concept}

\subsubsection{The window and perfect reconstruction}

TDAC cancels only if the analysis and synthesis windows satisfy the
Princen-Bradley conditions. For a window $w[n]$ of length $2N$, perfect
reconstruction requires
\begin{align}
  \label{eq:pr_symmetry}
  w[2N - 1 - n] &= w[n] && \text{(symmetry), and}\\
  \label{eq:pr_power}
  w[n]^2 + w[n + N]^2 &= 1 && \text{(power complementarity).}
\end{align}
CELT uses the \termdef{Vorbis window}, a power-complementary window built from a
sine of a sine:
\begin{equation}\label{eq:vorbis_window}
  w[n] = \sin\!\left[\frac{\pi}{2}\,
  \sin^2\!\left(\frac{\pi}{2L}\left(n + \tfrac{1}{2}\right)\right)\right],
  \quad n = 0,\dots,L-1,
\end{equation}
where $L$ is the overlap length. CELT does not window the full block: only the
first and last $L$ samples are tapered (with $L = 120$ at 48~kHz), and the
interior is passed through unwindowed. This \termdef{low-overlap}
design trades a
little frequency selectivity for substantially reduced algorithmic delay, which
is what lets Opus reach 5~ms latency in CELT mode.

\begin{implementation}{TDAC Reconstruction in the IMDCT}
  The synthesis overlap-add is the kernel that makes the whole transform
  invertible. After the IMDCT produces the time-domain block, the overlap region
  is mixed with the tail held over from the previous frame, applying the window
  to both edges so that \cref{eq:pr_power} cancels the aliasing:
  \begin{lstlisting}[language=Rust]
// Mirror both edges for TDAC, mixing with the previous frame's tail.
for i in 0..overlap / 2 {
    let (a, b) = (i, overlap - 1 - i);
    let (x1, x2) = (out[b], out[a]);
    let (w1, w2) = (window[i], window[overlap - 1 - i]);
    out[a] = w2 * x2 - w1 * x1;
    out[b] = w1 * x2 + w2 * x1;
}
  \end{lstlisting}
  The round trip (forward MDCT then inverse with overlap-add) is verified by a
  perfect-reconstruction test: a windowed block transformed and inverted returns
  the input to floating-point precision, with zero delay and unit gain.
\end{implementation}

\subsubsection{Short blocks for transients}

A long MDCT gives fine frequency resolution but coarse time resolution, which
smears a sharp transient (a drum hit, a consonant) backward in time as
\termdef{pre-echo}: quantization noise spread across the whole block becomes
audible before the attack that produced it. CELT's defense is to switch, for
transient frames, from one long MDCT to a series of $B$ short MDCTs (with $B =
2, 4, 8$), each covering a fraction of the frame. The encoder signals
this choice
with a transient flag (\cref{subsec:celt_encoder}); the short transforms'
coefficients are then interleaved so the band structure still applies. Time
resolution improves by a factor of $B$ at the cost of frequency resolution,
exactly the uncertainty trade-off, applied adaptively per frame.

\subsubsection{Computing the MDCT through an FFT}

Evaluated directly, \cref{eq:mdct} costs $O(N^2)$. In practice the MDCT
factors into three stages: a complex \emph{pre-rotation} that folds
the $2N$ real
inputs into $N/4$ complex values, an $N/4$-point complex FFT, and a complex
\emph{post-rotation} that unpacks the result. The FFT dominates the cost and is
isolated behind a seam in the implementation: the default build routes it
through an external planned-FFT backend, while a dependency-free
build falls back
to a direct $O(N^2)$ evaluation. The numerical and performance consequences of
this seam are treated in \cref{sec:optimizations}; here it suffices that the
transform is an $O(N\log N)$ operation in the production configuration.

\subsection{The Energy Envelope}\label{subsec:celt_energy}

CELT divides the MDCT spectrum into 21 \termdef{bands} whose widths follow the
ear's critical bands: narrow at low frequencies where the ear resolves pitch
finely, wide at high frequencies where it does
not~\cite{zwicker1999psychoacoustics}.
For each band, CELT codes two things separately: the band's \emph{energy} (the
$L^2$ norm of its coefficients) and its \emph{shape} (the unit-norm direction).
This separation is the ``Constrained-Energy'' in the codec's name.

\begin{concept}{Why Code Energy Separately}
  The audibility of a frequency band is governed largely by how much energy it
  contains, not by the precise arrangement of coefficients within it. By coding
  the energy envelope explicitly and protecting it with its own bits, CELT
  guarantees that the reconstructed signal has the right spectral balance even
  when the shape is coded coarsely. A codec that quantized raw
  coefficients would
  instead let energy drift as bits run short, producing audible
  spectral tilt and
  ``birdie'' artifacts. The cost of the guarantee is the overhead of
  the envelope
  itself, which CELT minimizes with aggressive prediction.
\end{concept}

Energy is coded in the $\log_2$ domain in three passes of decreasing
significance:

\begin{enumerate}[itemsep=3pt]
  \item \textbf{Coarse energy} quantizes each band's log-energy to roughly 6~dB
    resolution, predicted from neighbors and coded with a Laplace model.
  \item \textbf{Fine energy} refines each band with a number of raw bits set by
    the allocation, halving the error per bit.
  \item \textbf{Final bits} distribute any leftover budget, one bit at a time,
    to the bands that the allocation flagged as highest priority.
\end{enumerate}

\subsubsection{Coarse energy and two-dimensional prediction}

The coarse energy of a band is highly predictable from two neighbors: the same
band in the previous frame (time prediction) and the adjacent band in the same
frame (frequency prediction). CELT codes only the prediction residual, with a
Laplace probability model whose width depends on the frame size. The update
applies a prediction coefficient $\alpha$ in time and accumulates a running
frequency prediction with a leakage factor $\beta$:
\begin{equation}\label{eq:energy_pred}
  E_i \leftarrow \alpha\, E_i^{\text{prev}} + p_i + q_i,
  \qquad p_{i+1} = p_i + q_i(1 - \beta),
\end{equation}
where $q_i$ is the decoded residual for band $i$ and $p_i$ the frequency
predictor carried across bands. Intra frames (no time prediction) set $\alpha=0$
and use a larger $\beta$.

\begin{lstlisting}[language=Rust]
let (coef, beta) = if intra { (0.0, BETA_INTRA) }
                   else     { (PRED_COEF[lm], BETA_COEF[lm]) };
let mut prev = [0.0f32; 2];
for i in start..end {
    for (c, prev_c) in prev.iter_mut().enumerate().take(channels) {
        let qi = if budget - tell >= 15 {
            ec_laplace_decode(dec, /* model for band i, frame size lm */)
        } else { /* budget-starved fallback code */ };
        let old = state.old_ebands[c][i].max(-9.0);
        state.old_ebands[c][i] = coef * old + *prev_c + qi as f32;
        *prev_c += qi as f32 - beta * qi as f32;
    }
}
\end{lstlisting}

The Laplace coder (\S4.3.2.1) is a specialization of the range coder for a
two-sided geometric distribution: small residuals near zero are cheap, large
ones progressively more expensive, matching the statistics of inter-frame energy
deltas. When the remaining budget falls below the threshold to afford a Laplace
symbol, the coder drops to a fixed fallback code so the decode never stalls.

\subsection{Pyramid Vector Quantization}\label{subsec:pvq}

Once a band's energy is known, its shape is a unit vector in $N$ dimensions (for
a band of $N$ coefficients). CELT quantizes that direction with \termdef{pyramid
vector quantization} (PVQ)~\cite{fischer1986pyramid}, a structured vector
quantizer with a property that makes it ideal here: every codeword has exactly
the same norm, so quantizing the shape never disturbs the separately coded
energy.

\begin{concept}{The Pyramid Codebook}
  The PVQ codebook of dimension $N$ and \emph{pulse count} $K$ is the set of all
  integer vectors $\mathbf{y} \in \mathbb{Z}^N$ whose absolute components sum to
  $K$:
  \begin{equation}\label{eq:pvq_codebook}
    S(N, K) = \left\{\mathbf{y} \in \mathbb{Z}^N :
    \sum_{i=0}^{N-1} \abs{y_i} = K\right\}.
  \end{equation}
  Geometrically these are the lattice points on the surface of an
  $N$-dimensional
  $L^1$ ball, the ``pyramid.'' A band's shape is coded by finding the codeword
  $\mathbf{y}$ closest in direction to the unit shape, then transmitting which
  codeword it is. On decode, $\mathbf{y}$ is read back and normalized to unit
  norm, $\hat{\mathbf{x}} = \mathbf{y}/\norm{\mathbf{y}}$, and scaled by the
  band's energy. The pulse count $K$ controls the resolution: more pulses mean a
  denser codebook and a finer shape, and $K$ is exactly what the bit allocation
  (\cref{subsec:celt_allocation}) decides per band.
\end{concept}

\subsubsection{The pulse search}

Finding the best codeword is a search over $S(N,K)$, which would be enormous to
enumerate, so CELT uses a greedy algorithm that adds pulses one at a time, each
time placing a pulse in the position that most improves the correlation between
the candidate $\mathbf{y}$ and the target shape $\mathbf{x}$. After each pulse,
the quantity maximized is
\begin{equation}\label{eq:pvq_search}
  j^\star = \argmax_j \frac{(\mathbf{x}\cdot\mathbf{y} + x_j)^2}
  {\norm{\mathbf{y}}^2 + 2y_j + 1},
\end{equation}
the squared cosine similarity that would result from adding a pulse at position
$j$. The algorithm starts from a projection of $\mathbf{x}$ onto the pyramid to
place most of the pulses cheaply, then refines pulse by pulse. The inner loop of
\cref{eq:pvq_search} is the single hottest computation in CELT encoding and is
the subject of a dedicated SIMD kernel (\cref{subsec:simd_pvq}).

\subsection{CWRS: Enumerating the Codebook}\label{subsec:cwrs}

The PVQ search yields a codeword $\mathbf{y}$; the coder must
transmit \emph{which}
of the $V(N,K) = \abs{S(N,K)}$ codewords it is, as a uniform integer in
$[0, V(N,K))$. This is the \termdef{combinatorial} step: a bijection between
pulse vectors and indices, known in the Opus sources as CWRS (coding with
replacement and signs, \S4.3.4.2).

\begin{concept}{Counting Pulse Vectors}
  Let $V(N, K)$ be the number of length-$N$ integer vectors with $K$ signed
  pulses, and let $U(N, K)$ count those whose first nonzero entry is positive.
  These satisfy the recurrences
  \begin{align}
    \label{eq:cwrs_u}
    U(N, K) &= U(N-1, K) + U(N, K-1) + U(N-1, K-1), \\
    \label{eq:cwrs_v}
    V(N, K) &= U(N, K) + U(N, K+1),
  \end{align}
  with $U(0,0) = 1$. The three terms of \cref{eq:cwrs_u} correspond to the first
  position holding zero, holding a pulse with more to come, or holding a pulse
  that completes the budget. The index of a particular $\mathbf{y}$ is built by
  walking the positions and, at each, adding the count of all codewords that
  sort before the chosen one, the standard combinatorial-number-system ranking.
\end{concept}

Computing $U$ from \cref{eq:cwrs_u} on every band decode would cost $O(NK)$. The
implementation instead precomputes the $U$ rows into a flat table, so ranking
and unranking a codeword become $O(N)$ table lookups. The table has a little
over a thousand entries (1272 \code{u32}s) and exploits the symmetry
$U(n,k)=U(k,n)$ to stay compact. It is built by a \code{const fn} evaluated at
compile time, so the table is baked into the binary with no runtime
initialization and no allocator dependency (load-bearing for the
\code{no\_std} decoder core, \cref{sec:optimizations}):

\begin{lstlisting}[language=Rust]
// Row m's data starts at column m; PVQ_U_ROW[15] is the total length.
const fn compute_pvq_u_data() -> [u32; PVQ_U_DATA_LEN] {
    let mut uu = [[0u32; 177]; 15];   // U(a, b) for a <= 14, b <= 176
    uu[0][0] = 1;
    let mut a = 1;
    while a < 15 {
        let mut b = 1;
        while b < 177 {
            uu[a][b] = uu[a - 1][b].wrapping_add(uu[a][b - 1]).wrapping_add(uu[a - 1][b - 1]);
            b += 1;
        }
        a += 1;
    }
    let mut d = [0u32; PVQ_U_DATA_LEN];
    // ... pack row m of `uu` into d[PVQ_U_ROW[m]..], starting at column m ...
    d
}

static PVQ_U_DATA: [u32; PVQ_U_DATA_LEN] = compute_pvq_u_data();
\end{lstlisting}

The decoder reads the uniform index from the range coder, unranks it to recover
$\mathbf{y}$ through the table, and normalizes. Because the index is
uniform over
exactly $V(N,K)$ values, no bits are wasted: the shape costs precisely
$\log_2 V(N,K)$ bits, which is what the allocator budgeted for.

\subsection{Bit Allocation}\label{subsec:celt_allocation}

The energy envelope and the PVQ shapes both consume bits, and the number of
pulses $K$ available to each band is whatever the allocation leaves for it. The
allocation is the arbiter that converts a frame's total bit budget into a
per-band pulse count. Its defining constraint is subtle and important.

\begin{concept}{The Allocation Must Be Implicit}
  The decoder needs to know each band's pulse count \emph{before} it can decode
  that band's shape, but the pulse counts are never transmitted: there are no
  bits to spare for them. Instead, encoder and decoder run the \emph{same}
  deterministic allocation procedure over the same inputs (the total budget, the
  frame size, static tables, and a few coded parameters), and arrive at
  identical per-band allocations. This is why the bit budget is tracked to
  $1/8$-bit precision through \code{tell\_frac} (\cref{subsec:tell}): the
  allocation arithmetic must produce bit-identical results on both sides, or the
  decoder will read a band with the wrong pulse count and desynchronize from the
  range coder. The allocation is, in effect, a pseudo-random function of the
  budget that both parties compute independently.
\end{concept}

The procedure (\S4.3.3) works in units of $1/8$~bit (the constant
$\mathrm{BITRES}=3$). It begins from a family of static \termdef{allocation
vectors}, tabulated target bit counts per band for a range of quality levels,
and interpolates between two adjacent vectors to match the available budget. The
interpolation fraction is found by bisection: at each step it tentatively
allocates at the midpoint, sums the cost over all bands (clamped to per-band
caps, with bands below a floor skipped), and moves the bisection bound according
to whether the total overran the budget.

\begin{lstlisting}[language=Rust]
// Bisect the interpolation fraction in 1/64 steps to fit `total`.
let (mut lo, mut hi) = (0i32, 1 << ALLOC_STEPS);
for _ in 0..ALLOC_STEPS {
    let mid = (lo + hi) >> 1;
    let mut psum = 0i32;
    let mut done = false;
    for j in (start..end).rev() {
        let tmp = bits1[j] + ((mid * bits2[j]) >> ALLOC_STEPS);
        if tmp >= thresh[j] || done { done = true; psum += tmp.min(cap[j]); }
        else if tmp >= alloc_floor { psum += alloc_floor; }
    }
    if psum > total { hi = mid; } else { lo = mid; }
}
\end{lstlisting}

Once the per-band bit targets are fixed, each band's allocation is split between
\emph{fine-energy} bits (raw bits that refine the envelope) and \emph{shape}
bits (the PVQ codeword). The conversion between a bit budget and a pulse count
$K$ goes through a precomputed \termdef{pulse cache}: \code{bits2pulses} binary-
searches the cache for the largest $K$ whose codeword fits the budget, and
\code{pulses2bits} reports the exact cost $\lceil\log_2 V(N,K)\rceil$ of a given
$K$. Surplus bits left after the split, the \emph{balance}, are carried forward
to subsequent bands so nothing is wasted. The allocation also decides, from the
budget, the boundary band for intensity stereo and whether to use dual stereo
(\cref{subsec:celt_bands}), and it may explicitly \emph{skip} high
bands when the
budget cannot support them, coding the skip decisions as run-length flags.

\subsection{Band Coding: Splits, Folding, and Stereo}\label{subsec:celt_bands}

With energy and allocation settled, the band loop codes each band's shape. A
band of $N$ coefficients is rarely quantized as one PVQ vector; instead it is
recursively halved, with the split coded by an angle, until the pieces are small
enough to quantize directly. Several refinements layer on top: spreading to
decorrelate noisy bands, folding to fill empty ones, anti-collapse to repair
transient artifacts, and stereo coupling.

\subsubsection{Recursive splitting by angle}

When a band is split in two, the energy is divided between the halves
by an angle
$\theta$: the lower half receives $\cos\theta$ of the amplitude and the upper
$\sin\theta$. Quantizing $\theta$ (rather than two independent energies) keeps
the total band energy exactly fixed, preserving the energy constraint
through the
recursion. The number of quantization levels for $\theta$ depends on the band's
bit allocation, and the angle is coded with one of three probability models
chosen by context (\S4.3.4.1): a \emph{uniform} PDF for ordinary time splits, a
\emph{triangular} PDF that biases toward an even split for mono bands, and a
\emph{step} PDF for the first stereo split. The recursion terminates when a sub-
band is small enough, at which point it is coded by a single PVQ vector
(\cref{subsec:pvq}).

\subsubsection{Spreading rotation}

Concentrating $K$ pulses in $N$ positions produces a sparse, spiky vector, which
in a band that should sound noise-like is audible as tonal ``birdies.'' CELT
counters this by applying a \termdef{spreading rotation}: a sequence of Givens
rotations that smear the pulses' energy across the band without changing its
norm. The amount of spreading (none, light, normal, or aggressive) is chosen by
the encoder's tonality analysis (\cref{subsec:celt_encoder}) and signaled once
per frame.

\subsubsection{Folding and noise fill}

When the allocation gives a band zero pulses, leaving it empty would create
silent spectral holes. Instead CELT fills the band by \termdef{folding}: copying
already-decoded coefficients from lower frequencies into the empty band, then
renormalizing to the band's coded energy. When no suitable lower content exists,
the band is filled with pseudo-random noise from a small linear congruential
generator. The generator and the fold share the same role, supplying a plausible
shape at the correct energy for bands too cheap to code explicitly.

\begin{lstlisting}[language=Rust]
// Linear congruential generator used for noise fill (Numerical Recipes constants).
fn celt_lcg_rand(seed: u32) -> u32 {
    seed.wrapping_mul(1_664_525).wrapping_add(1_013_904_223)
}
\end{lstlisting}

\subsubsection{Collapse masks and anti-collapse}

In a transient frame coded as $B$ short blocks (\cref{subsec:mdct}), a band's
pulses may all land in a single short block, leaving the others at zero, a
\termdef{collapse}. A collapsed block is silent for part of the frame, which can
sound like a dropout. CELT tracks a \termdef{collapse mask}, one bit per short
block per band, recording which blocks received energy. The \termdef{anti-
collapse} step then injects a controlled amount of noise into the collapsed
blocks, scaled to the minimum of the neighboring bands' energies, so every short
block carries at least a floor of energy. The presence of anti-collapse is
controlled by a single coded bit per transient frame.

\subsubsection{Stereo coupling}

For stereo input, CELT codes each band as a \emph{mid} (sum) and \emph{side}
(difference) pair rather than left and right, because the mid/side
representation concentrates energy in the mid channel and codes the stereo image
efficiently. Three sub-modes apply, chosen per band by the allocation:
\begin{itemize}[itemsep=2pt]
  \item \textbf{Mid/side with angle:} a stereo angle $\theta$ apportions energy
    between mid and side, coded like the mono split above.
  \item \textbf{Intensity stereo:} above a boundary band set by the budget, only
    the mid is coded and the side is reconstructed by copying the
    mid, discarding
    stereo detail where it is least audible (at high frequencies, low bitrates).
  \item \textbf{Dual stereo:} the two channels are coded fully
    independently, used
    when the channels are sufficiently decorrelated to justify the bits.
\end{itemize}
A special case handles bands of $N=2$, where a single angle suffices to code the
stereo relationship directly.

\subsubsection{Time-frequency reshaping}

Within a frame, individual bands can opt for finer time resolution (at the
expense of frequency resolution) or vice versa, independently of the global
short-block decision. This is implemented with in-place Haar and Hadamard
transforms that redistribute a band's coefficients between time and frequency
before quantization. A per-frame \code{tf\_select} flag and per-band flags
encode the choice, which the encoder optimizes by the analysis below.

\subsection{The Encoder's Analysis}\label{subsec:celt_encoder}

Everything to this point is shared by encoder and decoder; the decoder simply
reads the parameters and follows. The encoder must \emph{decide} those
parameters, and those decisions, not the bitstream format, are where encoders
differ in quality. CELT's encoder runs several analyses per frame.

\textbf{Transient detection.} The encoder decides whether to use short blocks by
running the input through a high-pass filter and then forward and
backward energy
followers; a sharp rise in the high-passed energy relative to its smoothed
envelope marks a transient. The high-pass is a second-order section, and its
output drives the masking metric:

\begin{lstlisting}[language=Rust]
// High-pass filter: (1 - 2z^-1 + z^-2) / (1 - z^-1 + 0.5 z^-2).
let (mut mem0, mut mem1) = (0.0f32, 0.0f32);
for (t, &x) in tmp.iter_mut().zip(input.iter()) {
    let y = mem0 + x;
    let mem00 = mem0;
    mem0 = mem0 - x + 0.5 * mem1;
    mem1 = x - mem00;
    *t = y;
}
\end{lstlisting}

\textbf{Time-frequency analysis.} For the bands, the encoder chooses
time-versus-
frequency resolution by an $L^1$-metric minimization solved with a small Viterbi
search (\code{tf\_analysis} and \code{tf\_select}), trading the cost of changing
resolution against the coding gain it yields band by band.

\textbf{Spreading decision.} A tonality measure over the spectrum selects the
spreading amount (\cref{subsec:celt_bands}): tonal frames get little or no
spreading, noisy frames get aggressive spreading.

\textbf{Dynamic allocation and trim.} Two analyses shape the
allocation before it
runs. \termdef{Dynamic allocation} (\code{dynalloc\_analysis}) computes per-band
importance boosts from a masking model, steering extra bits to perceptually
important bands. The \termdef{allocation trim} (\code{alloc\_trim\_analysis})
sets a single global parameter that tilts the budget toward low or high
frequencies based on spectral tilt and transientness.

\textbf{The pitch pre-filter.} Before the MDCT, the encoder optionally applies a
comb \termdef{pre-filter} that whitens the signal's harmonic
structure, the exact
inverse of the comb post-filter the decoder applies after synthesis. Whitening
the harmonics concentrates energy and improves coding efficiency on strongly
tonal material. The pitch period and gain that parameterize the filter come from
a pitch analysis (downsampling, decimated cross-correlation, and octave-error
rejection) that mirrors the technique used for packet-loss concealment
(\cref{subsec:plc}).

\section{SILK: Linear-Predictive Coding}\label{sec:silk}

SILK (\S4.2) is the speech half of Opus. Where CELT works in the frequency
domain on arbitrary signals, SILK works in the time domain on a \emph{model} of
how speech is produced, and codes only what the model fails to predict. This
model-based approach is what lets SILK code intelligible speech at a few
kilobits per second, far below what a general transform coder
achieves. SILK runs
at one of three internal rates (8, 12, or 16~kHz, for narrowband,
  mediumband, and
wideband), in frames of 10, 20, 40, or 60~ms divided into 2 or 4 subframes.

This section develops the model bottom-up: the source-filter idea
(\cref{subsec:silk_model}), the short-term predictor and how it is estimated and
represented (\cref{subsec:silk_lpc,subsec:silk_nlsf}), the long-term predictor
that captures pitch (\cref{subsec:silk_ltp}), the gains and excitation coding
(\cref{subsec:silk_gains,subsec:silk_excitation}), the noise-shaping
quantizer at
the encoder's heart (\cref{subsec:silk_nsq}), and the supporting machinery of
stereo, resampling, and control
(\cref{subsec:silk_stereo,subsec:silk_resampler,subsec:silk_control}).

\subsection{The Source-Filter Model}\label{subsec:silk_model}

\begin{concept}{Speech as Source and Filter}
  Human speech is produced by an excitation source (the vocal folds, for voiced
  sounds, or turbulent airflow, for unvoiced ones) passing through a filter (the
  vocal tract, whose shape determines the formants we hear as vowels). Linear-
  predictive coding~\cite{makhoul1975linear,rabiner1978digital} models this
  directly: the vocal tract becomes an all-pole filter $1/A(z)$, and the
  excitation becomes a sparse residual signal. The model predicts each sample
  from a weighted sum of recent samples,
  \begin{equation}\label{eq:lpc_predict}
    \hat{s}[n] = -\sum_{k=1}^{p} a_k\, s[n-k],
  \end{equation}
  and codes only the residual $e[n] = s[n] - \hat{s}[n]$ that the prediction
  misses. Because the residual carries far less energy and structure than the
  original, it is cheap to code. SILK adds a second predictor for the residual's
  own periodicity (pitch) and a noise-shaping loop that hides quantization noise
  under the speech spectrum.
\end{concept}

The frame is the unit at which SILK transmits its model; the subframe
is the unit
at which the model is applied. Each frame carries one set of short-term filter
coefficients (interpolated across subframes for 20~ms frames), and per-subframe
gains and pitch parameters. The decoder reconstructs each frame by running the
excitation through the long-term and short-term synthesis filters, then
resampling to the output rate.

\subsection{Short-Term Prediction}\label{subsec:silk_lpc}

The short-term predictor of \cref{eq:lpc_predict} is a $p$-tap linear filter
($p=10$ for narrowband and mediumband, $p=16$ for wideband). It comes in two
forms that are exact inverses: the \termdef{analysis} filter $A(z)$, which the
encoder applies to produce the residual, and the \termdef{synthesis} filter
$1/A(z)$, which the decoder applies to reconstruct the signal. The
implementation runs this in \Q{12} fixed-point on \code{i16} coefficients and
samples (SILK's convention predicts with a \emph{plus} sign, so the analysis
  filter subtracts the prediction from the input rather than adding it, unlike
\cref{eq:lpc_predict}'s $A(z)$ sign). The taps are reversed once up front so
the per-sample dot product is a bounds-check-free forward window:

\begin{lstlisting}[language=Rust]
// MA whitening filter (Q12 taps); the first `order` outputs are zeroed.
pub(crate) fn lpc_analysis_filter(out: &mut [i16], input: &[i16], b: &[i16]) {
    let d = b.len();
    let mut b_rev = [0i16; 16];
    for (j, &tap) in b.iter().enumerate() {
        b_rev[d - 1 - j] = tap;
    }
    for ix in d..input.len() {
        let win = &input[ix - d..ix];
        let mut pred_q12 = 0i32;
        for (&w, &tap) in win.iter().zip(b_rev[..d].iter()) {
            pred_q12 = smlabb(pred_q12, i32::from(w), i32::from(tap));
        }
        let e_q12 = (i32::from(input[ix]) << 12) - pred_q12;   // e[n] = x[n] - pred
        out[ix] = rshift_round(e_q12, 12).clamp(i32::from(i16::MIN), i32::from(i16::MAX)) as i16;
    }
    out[..d].fill(0);
}
\end{lstlisting}

The synthesis filter run by the decoder is the same recurrence with the sign
and the roles of input/output swapped: it accumulates the prediction from
\emph{already-reconstructed} output samples and adds it back to the incoming
residual, so past decisions feed forward exactly as they did during encoding.

\subsubsection{Estimating the coefficients}

The decoder receives the coefficients; the encoder must estimate them from the
signal. Two classical methods apply. \termdef{Levinson-Durbin}
recursion~\cite{levinson1947wiener} solves the normal equations relating the
filter coefficients to the signal's autocorrelation, in $O(p^2)$ time, producing
the minimum mean-square prediction error. \termdef{Burg's
method}~\cite{burg1975maximum} instead estimates reflection
coefficients directly
from the signal by minimizing the sum of forward and backward prediction errors,
which behaves better on the short analysis windows speech coding uses. The SILK
encoder uses a modified Burg method, and caps the resulting filter's prediction
gain to keep it numerically well-conditioned.

\begin{implementation}{Levinson-Durbin Recursion}
  The recursion builds the order-$p$ predictor from the order-$(p-1)$ predictor
  by computing a reflection coefficient $k_m$ that drives the next prediction
  error to its minimum, updating all coefficients in place:
  \begin{lstlisting}[language=Rust]
pub fn levinson_durbin(autocorr: &[f64], order: usize) -> Vec<f64> {
    let mut a = vec![0.0; order + 1];
    let mut err = autocorr[0];
    for m in 1..=order {
        let mut k = autocorr[m];
        for i in 1..m { k += a[i] * autocorr[m - i]; }
        k = -k / err;                       // reflection coefficient
        a[m] = k;
        for i in 1..=m / 2 {                // symmetric in-place update
            let tmp = a[i] + k * a[m - i];
            a[m - i] += k * a[i];
            a[i] = tmp;
        }
        err *= 1.0 - k * k;                 // error shrinks each order
    }
    a
}
  \end{lstlisting}
  The factor $1 - k_m^2$ is non-negative exactly when $\abs{k_m} \le 1$, the
  condition for filter stability, which is why reflection coefficients are a
  convenient parameterization.
\end{implementation}

\subsection{Line Spectral Frequencies}\label{subsec:silk_nlsf}

The encoder cannot transmit LPC coefficients directly: they are extremely
sensitive to quantization, and a small rounding error can push the all-pole
synthesis filter outside the unit circle, making it unstable and the decoded
signal explode. SILK transmits the filter in a different representation, the
\termdef{line spectral frequencies} (LSFs), which is built for quantization.

\begin{concept}{Why Line Spectral Frequencies}
  The order-$p$ polynomial $A(z)$ is decomposed into a symmetric and an
  antisymmetric polynomial,
  \begin{align}
    P(z) &= A(z) + z^{-(p+1)} A(z^{-1}), \\
    Q(z) &= A(z) - z^{-(p+1)} A(z^{-1}),
  \end{align}
  whose roots, the LSFs, all lie \emph{on} the unit circle and
  \emph{interlace}~\cite{itakura1975line}. This representation has three
  properties that make it ideal for coding. First, stability is trivially
  checked and enforced: as long as the LSFs stay ordered, the reconstructed
  filter is stable. Second, the LSFs are localized in frequency, so each one
  affects the spectrum mainly near its own frequency, and a
  quantization error in
  one LSF stays local rather than corrupting the whole spectrum. Third, they
  interpolate gracefully, so SILK can transmit one set per frame and interpolate
  across subframes. SILK uses \termdef{normalized} LSFs (NLSFs), scaled so the
  frequency range maps to $[0, 1)$ and stored as \Q{15} integers.
\end{concept}

\subsubsection{The two-stage vector quantizer}

NLSFs are coded with a two-stage vector quantizer (\S4.2.7.5). The first stage
selects one of 32 codebook vectors that approximates the whole NLSF vector,
coded with an entropy model. The second stage codes the per-coefficient residual
between the true NLSFs and the first-stage approximation, using a predictive
scheme (each residual partly predicted from the previous) and an
arithmetic-coded
index. The residual quantization is weighted by \termdef{Laroia
weights}~\cite{laroia1991robust}, which approximate each
coefficient's perceptual
sensitivity (inversely related to its spacing from its neighbors), so bits are
spent where the ear notices. The encoder searches this quantizer with a delayed-
decision trellis, keeping several candidate paths alive to find a low-rate,
low-distortion combination.

\subsubsection{Reconstruction and stability enforcement}

The decoder reverses the process: dequantize the residual, add the first-stage
vector, enforce a minimum spacing between adjacent NLSFs to guarantee stability,
and convert back to LPC coefficients. The conversion (\code{nlsf2a}) builds the
$P$ and $Q$ polynomials from the cosines of the NLSFs and combines them. Even
after stabilization, the resulting filter is checked for stability, and if it
fails, bandwidth expansion (scaling the coefficients toward zero) is applied
until it passes:

\begin{lstlisting}[language=Rust]
// After converting NLSFs to LPC, force stability by bandwidth expansion.
while lpc_inverse_pred_gain(&a_q12[..d]) == 0 && i < MAX_LPC_STABILIZE_ITERATIONS {
    bwexpander_32(&mut a32_qa1[..d], 65536 - (2 << i));  // shrink toward zero
    /* re-quantize a_q12 from a32_qa1 and re-test */
    i += 1;
}
\end{lstlisting}

For 20~ms frames, the decoder interpolates between the previous
frame's NLSFs and
the current frame's, using a 2-bit interpolation factor, so the filter evolves
smoothly across the frame's first and second halves rather than jumping at the
midpoint.

\subsection{Long-Term Prediction}\label{subsec:silk_ltp}

Short-term prediction removes the spectral envelope, but the residual of voiced
speech is still strongly periodic at the pitch period: the vocal folds produce a
pulse train, and successive pitch periods resemble one another. SILK
removes this
redundancy with a second predictor, the \termdef{long-term predictor} (LTP),
which predicts the current residual from the residual one pitch period earlier.

\begin{concept}{Pitch Prediction with a Five-Tap Filter}
  For a voiced subframe with pitch lag $T$ (in samples), the LTP predicts the
  residual $e[n]$ from a five-tap filter applied around $e[n-T]$:
  \begin{equation}\label{eq:ltp}
    \hat{e}[n] = \sum_{j=-2}^{2} b_j\, e[n - T + j].
  \end{equation}
  The five taps $b_j$ let the predictor track a pitch period that does not fall
  exactly on a sample boundary and that changes shape slowly from period to
  period. What remains after long-term prediction, the
  doubly-predicted residual,
  is the sparse excitation that SILK actually codes. The pitch lag $T$ and the
  taps $b_j$ are themselves transmitted: the lag as a per-frame absolute value
  plus per-subframe offsets from a contour codebook, and the taps as an index
  into one of three gain codebooks.
\end{concept}

The lag is decoded as a base lag plus a per-subframe contour offset, clamped to
the valid pitch range for the internal rate:

\begin{lstlisting}[language=Rust]
let (min_lag, max_lag) = (PE_MIN_LAG_MS * fs_khz, PE_MAX_LAG_MS * fs_khz);
let lag = min_lag + i32::from(lag_index);
for (k, out) in pitch_lags.iter_mut().enumerate().take(nb_subfr) {
    *out = (lag + contour_offset(k)).clamp(min_lag, max_lag);  // per-subframe lag
}
\end{lstlisting}

The five-tap gains come from one of three codebooks of 8, 16, or 32 vectors,
selected by a per-frame \emph{periodicity index} that trades codebook size
against the bits to address it: highly periodic speech justifies the
larger, more
precise codebook. The encoder chooses the lag with a three-stage pitch search
(\cref{subsec:silk_control}) and the gains with a
rate-distortion-weighted vector
quantizer.

\subsection{Subframe Gains}\label{subsec:silk_gains}

Each subframe carries a gain that scales the excitation before it enters the
synthesis filters, tracking the signal's changing loudness. Gains are quantized
in the $\log_2$ domain across 64 levels, so the quantization step is a constant
ratio rather than a constant difference, which matches the multiplicative way
loudness varies. After the first subframe, gains are delta-coded against the
previous subframe, with a wider ``double step'' available for rapid changes. The
decoder converts the log-domain index back to a linear \Q{16} gain through a
piecewise-parabolic $\log_2$-to-linear approximation:

\begin{lstlisting}[language=Rust]
for k in 0..nb_subfr {
    if k == 0 && !conditional {
        // The index may not drop more than 16 steps (~21.8 dB).
        *prev_ind = i32::from(ind[k]).max(i32::from(*prev_ind) - 16) as i8;
    } else {
        let ind_tmp = i32::from(ind[k]) + MIN_DELTA_GAIN_QUANT;   // delta-coded
        // Accumulate deltas, with a wider double step above a threshold.
        let double_step_threshold =
            2 * MAX_DELTA_GAIN_QUANT - N_LEVELS_QGAIN + i32::from(*prev_ind);
        *prev_ind = if ind_tmp > double_step_threshold {
            i32::from(*prev_ind) + ((ind_tmp << 1) - double_step_threshold)
        } else {
            i32::from(*prev_ind) + ind_tmp
        } as i8;
    }
    *prev_ind = i32::from(*prev_ind).clamp(0, N_LEVELS_QGAIN - 1) as i8;
    // Scale into the log domain and back to linear Q16.
    gain_q16[k] = log2lin((smulwb(INV_SCALE_Q16, i32::from(*prev_ind)) + OFFSET).min(3967));
}
\end{lstlisting}

\subsection{The Excitation and the Shell Coder}\label{subsec:silk_excitation}

After both predictors and the gain, what remains to code is the
\termdef{excitation}: the quantized residual, a mostly-small integer per sample.
SILK codes it in 16-sample blocks with a \termdef{shell coder} (\S4.2.7.8), a
recursive scheme that codes the \emph{distribution} of pulses within a block
rather than each sample independently.

\begin{concept}{Recursive Pulse Splitting}
  Given the total number of pulses in a 16-sample block, the shell coder splits
  that total down a binary tree: the 16-sample count divides into two 8-sample
  counts, each into two 4-sample counts, and so on down to individual samples.
  At each node only the \emph{split} is coded, that is, how the parent's pulse
  count divides between its two children, using a probability model
  conditioned on
  the parent count. This is efficient because the total pulse count already
  constrains the splits: a block with two pulses has far fewer ways
  to distribute
  them than a block with ten, and the conditional models capture exactly that.
  Pulse magnitudes beyond what the tree codes directly are sent as additional
  least-significant-bit planes, and each non-zero sample's sign is coded
  afterward with a probability that depends on the signal type, the quantization
  offset, and the block's pulse count.
\end{concept}

\begin{lstlisting}[language=Rust]
// One node of the shell tree: split a parent count into two child counts.
fn decode_split(dec: &mut RangeDecoder, p: i32, shell_table: &[u8]) -> (i32, i32) {
    if p > 0 {
        let off = SHELL_CODE_TABLE_OFFSETS[p as usize] as usize;
        let child1 = dec.decode_icdf(&shell_table[off..], 8) as i32;
        (child1, p - child1)            // the other child is the remainder
    } else {
        (0, 0)
    }
}
\end{lstlisting}

A per-frame \termdef{rate level} selects which family of shell probability
tables to use, letting the excitation coder adapt its models to the bitrate.

\subsection{The Noise-Shaping Quantizer}\label{subsec:silk_nsq}

The noise-shaping quantizer (NSQ) is the central kernel of the SILK encoder, the
stage that actually decides the excitation values. A naive quantizer would round
the residual to the nearest level independently per sample; the NSQ does
something more sophisticated, choosing each excitation sample to shape the
spectrum of the resulting quantization \emph{noise} so that it hides beneath the
speech.

\begin{concept}{Shaping Quantization Noise}
  Quantization injects error, and that error has a spectrum. If the error
  spectrum is flat (white), it is audible wherever the speech is quiet. The
  insight of noise shaping is to spectrally color the error so it follows the
  speech spectrum, concentrating noise at frequencies where the speech is loud
  enough to mask it~\cite{rabiner1978digital}. The NSQ achieves this by running
  the quantization inside a feedback loop: the error from past samples is
  filtered by a \termdef{noise-shaping filter} and fed back into the current
  prediction, so the quantizer is steered to make errors whose
  spectrum matches a
  perceptual target. Because the loop reconstructs exactly the signal
  the decoder
  will produce, the encoder always knows the true reconstruction state, which is
  what allows the per-sample decisions to be optimal rather than open-loop.
\end{concept}

For each sample the NSQ forms the short-term and long-term predictions, adds the
shaping feedback, quantizes the result by a rate-distortion rule (minimizing
  $\abs{\text{error}} + \lambda \cdot \text{bits}$, where $\lambda$
  comes from the
target coding quality), emits the excitation value, and updates its filter
states. The prediction and feedback dot products are computed in fixed point,
bit-identically to the reference, so the encoder's reconstruction matches the
decoder's to the sample:

\begin{lstlisting}[language=Rust]
// The LPC prediction (Q10). `coef_rev` is reversed once per subframe so this
// descending-history access becomes a forward windowed dot with no bounds
// checks in the per-sample loop.
fn short_prediction(buf: &[i32], base: usize, coef_rev: &[i16]) -> i32 {
    let order = coef_rev.len();
    let w = &buf[base + 1 - order..=base];
    let mut out = (order >> 1) as i32;
    for (&wj, &cj) in w.iter().zip(coef_rev.iter()) {
        out = smlawb(out, wj, i32::from(cj));   // scalar SMLAWB chain
    }
    out
}
\end{lstlisting}

This one is deliberately scalar rather than routed through the SIMD dot
products of \cref{subsec:simd_dot}: it runs once per sample over at most 16
taps, and a SIMD call's fixed overhead (the call itself, and folding the
lanes back to a scalar) does not amortize over so few taps. The scalar/SIMD
line is drawn per call site by measured cost, not by a blanket rule.

\subsubsection{Noise-shaping analysis}

The filters the NSQ applies are designed once per frame by the \termdef{noise-
shaping analysis}. It examines the speech and produces four shaping controls: an
auto-regressive (AR) filter derived from the LPC envelope and bandwidth-expanded
to set the overall noise shape; a spectral \termdef{tilt} that emphasizes or de-
emphasizes low frequencies; a low-frequency shaper; and, for voiced frames, a
\termdef{harmonic shaper} that shapes noise around the pitch
harmonics. These are
driven by the target SNR and the voice-activity analysis so that the noise
follows the masking threshold as the signal changes.

\subsection{Stereo Prediction}\label{subsec:silk_stereo}

Stereo SILK codes a mid channel (the average of left and right) and a predicted
side channel (their difference), because the side is usually well predicted from
the mid and so costs few bits. The encoder estimates per-frame
prediction weights
that reconstruct the side from a three-tap smoothed mid plus a one-tap lagged
mid; the weights are quantized, coded, and interpolated over the frame's first
8~ms to avoid discontinuities. When the side is sufficiently predictable, a
\termdef{mid-only} flag drops it entirely. The decoder applies the inverse:

\begin{lstlisting}[language=Rust]
// Reconstruct the side channel from the mid using interpolated predictor weights,
// then convert mid/side back to left/right.
for n in 0..frame_length {
    let sum  = i32::from(x1[n + 1]) + i32::from(x2[n + 1]);  // L = mid + side
    let diff = i32::from(x1[n + 1]) - i32::from(x2[n + 1]);  // R = mid - side
    x1[n + 1] = sum.clamp(i16::MIN as i32, i16::MAX as i32) as i16;
    x2[n + 1] = diff.clamp(i16::MIN as i32, i16::MAX as i32) as i16;
}
\end{lstlisting}

\subsection{Resampling}\label{subsec:silk_resampler}

SILK runs at 8, 12, or 16~kHz internally, but the caller asks for 8, 12, 16, 24,
or 48~kHz output, so the decoder ends with a resampling stage (\S4.2.8). The
method depends on the ratio: equal rates pass through, a factor-of-two change
uses a cascade of second-order allpass sections (cheaper and lower-delay than a
long FIR for exact 2$\times$), and other ratios combine the 2$\times$
allpass with
a fractional-delay FIR interpolation. The encoder's front-end runs the inverse,
decimating 48~kHz input down to the internal rate; its half-band and 2/3-band
decimators are pinned bit-exactly against the reference, since the resampled
signal is what the rest of the encoder analyzes.

\subsection{Voice Activity, Pitch, and Rate Control}\label{subsec:silk_control}

Three encoder analyses decide the parameters the rest of SILK consumes.

\textbf{Voice activity detection.} The VAD splits each frame into four frequency
bands with a tree of allpass half-band filters, tracks the noise floor in each
band, and maps the per-band signal-to-noise ratios through a sigmoid to a
speech-activity probability (a \Q{8} value). This probability, along with a
spectral-tilt measure, drives the noise-shaping analysis, the gain control, and
the pitch threshold, so that the codec spends its effort where there is speech.

\textbf{Pitch analysis.} The pitch lag is found by a three-stage search of
increasing precision: a coarse cross-correlation at 4~kHz proposes candidate
lags, a refinement at 8~kHz searches a small contour codebook around them, and a
final search at the internal rate locks in the per-subframe lags. The
search also
classifies the frame as voiced or unvoiced, which selects whether the long-term
predictor is used at all.

\textbf{Rate control.} A target bitrate is translated into a target coding SNR
through per-bandwidth tables, with an adjustment for the shorter 10~ms frames.
The coding SNR then sets the rate-distortion trade-off $\lambda$ in the NSQ and
the aggressiveness of the noise shaping. When a coded frame would exceed a hard
byte budget, the gains are scaled coarser and the frame re-encoded
until it fits,
which is how SILK honors a bitrate cap and how the hybrid mode keeps room for
CELT's high band (\cref{subsec:hybrid}).

\subsection{Packet-Loss Concealment in SILK}\label{subsec:silk_plc}

When a SILK frame is lost, the decoder conceals it by extrapolating
the last good
frame's model: it continues the pitch excitation with progressively stronger
attenuation, drifts the pitch lag, and fades toward the spectral envelope of the
preceding frames, mixing in comfort noise estimated from the recent signal. The
attenuation and pitch-drift constants are tuned so a short loss is inaudible and
a long one decays gracefully to silence rather than repeating a buzz.
This is the
SILK-specific counterpart to the general concealment logic in
\cref{subsec:plc}.

\section{The Opus Layer}\label{sec:opus_layer}

SILK and CELT are complete coders on their own. The Opus layer (\S4 for the
decoder, \S5 for the encoder) is what makes them one codec: it dispatches each
frame to the right coder, runs both together in hybrid mode, smooths the
transitions when the mode changes between frames, and implements the resilience
features (concealment, forward error correction, discontinuous transmission)
that operate above either coder. This layer holds the cross-frame state that
neither coder can see on its own.

\subsection{Mode Dispatch}\label{subsec:dispatch}

The TOC byte names the mode; the decoder dispatches accordingly. A frame's bytes
become a single range decoder, and the SILK half (if any) decodes first,
followed by the CELT half (if any), both drawing from the same range decoder so
the bitstream is one continuous arithmetic-coded message. After both halves, the
internal output is resampled to the API rate and any cross-frame smoothing is
applied. The skeleton of the per-frame decode is:

\begin{lstlisting}[language=Rust]
let mut dec = RangeDecoder::new(data);
if mode != Mode::CeltOnly {
    // SILK half: low band (SILK-only) or 0..8 kHz (hybrid), at 8/12/16 kHz.
    self.silk.decode(&mut dec, &ctl, /* first call */, &mut pcm_silk);
    for (p, &s) in pcm.iter_mut().zip(pcm_silk.iter()) {
        *p += f32::from(s) / 32768.0;
    }
}
// ... optional redundancy handling ...
if mode != Mode::SilkOnly {
    // CELT half: full band (CELT-only) or bands 17..end (hybrid).
    let start_band = if mode == Mode::CeltOnly { 0 } else { 17 };
    pcm = self.celt.decode_frame(&mut dec, len, frame_size,
                                 self.stream_channels, start_band, celt_end);
}
self.final_range = dec.range_size() ^ redundant_rng;
\end{lstlisting}

\subsection{Hybrid Mode and the Crossover}\label{subsec:hybrid}

Hybrid mode is the reason SILK and CELT share a range coder. In a hybrid frame,
SILK codes the signal up to roughly 8~kHz and CELT codes everything above, and
the two contributions are summed (SILK's after resampling its 16~kHz output to
48~kHz) to reconstruct the full-band signal.

\begin{concept}{Splitting the Spectrum Between Two Coders}
  CELT's band structure assigns a fixed range of MDCT bands to each frequency
  region. In hybrid mode, CELT simply starts at band 17 (about 8~kHz) instead of
  band 0, leaving the bands below for SILK. The crossover is not a filter but a
  bookkeeping boundary: SILK is given the low band as its full-band wideband
  signal, and CELT is told to begin coding at band 17. Because both write into
  one range coder, the decoder reads the SILK parameters, then the CELT
  parameters, with no marker between them; the byte where SILK ends and CELT
  begins is wherever the SILK decode happens to leave the range coder. The
  highest CELT band depends on the audio bandwidth: super-wideband stops at band
  19, fullband at band 21.
\end{concept}

The encoder side must budget the shared bytes so that CELT always has room for
its high band. It computes a SILK bitrate from a rate table, caps the SILK low
band to a hard byte budget, and lets CELT fill the remainder, all into one range
encoder:

\begin{lstlisting}[language=Rust]
let celt_floor = (celt_end - 17) * 3 + 3;          // min bytes for the high band
let silk_cap = nb_bytes.saturating_sub(celt_floor).max(8);
let mut enc = RangeEncoder::new(nb_bytes);
silk.encode_into(&mut enc, &internal, Some((silk_cap * 8) as i32));  // capped
self.celt.encode_hybrid_into(&mut enc, pcm, nb_bytes, celt_end);     // fills rest
\end{lstlisting}

\subsection{Mode Transitions and Redundancy}\label{subsec:redundancy}

Because the mode can change from one packet to the next, the decoder must bridge
discontinuities when, say, a SILK-only packet is followed by a CELT-only packet.
The coders carry independent internal state (filter memories, overlap buffers),
and switching cold would click. Opus addresses this with two mechanisms: state
resets at the boundary, and an optional \termdef{redundancy} frame.

\begin{concept}{Redundancy Frames at Transitions}
  A redundancy frame is a short (5~ms) CELT frame embedded in a SILK or hybrid
  packet, used to fill the gap during a mode switch. When switching from CELT to
  SILK, the redundant CELT frame covers the start of the SILK packet, which has
  no valid CELT history; when switching from SILK to CELT, it covers
  the end. The
  decoder cross-fades between the redundant frame and the primary decode over a
  2.5~ms overlap using a squared Vorbis window, so the energy sums to
  unity across
  the splice and no discontinuity is heard. The presence and direction of
  redundancy are coded with a few bits at the boundary of the SILK
  and CELT data.
\end{concept}

The crossfade is the same power-complementary mix used throughout Opus, applied
sample by sample over the overlap:

\begin{lstlisting}[language=Rust]
fn smooth_fade(in1: &[f32], in2: &[f32], out: &mut [f32],
               overlap: usize, channels: usize, inc: usize) {
    for c in 0..channels {
        for i in 0..overlap {
            let w = WINDOW120[i * inc] * WINDOW120[i * inc];   // squared window
            out[i * channels + c] =
                w * in2[i * channels + c] + (1.0 - w) * in1[i * channels + c];
        }
    }
}
\end{lstlisting}

\subsection{Packet-Loss Concealment}\label{subsec:plc}

When a packet does not arrive, the decoder must produce a plausible frame of
audio rather than a gap. This is \termdef{packet-loss concealment} (PLC). Opus
conceals in whichever mode was last active: the CELT concealer extrapolates the
last frame's pitch period with decaying gain, and the SILK concealer continues
its excitation model (\cref{subsec:silk_plc}). A loss is signaled to the decoder
by handing it an empty or single-byte packet, which it routes to the concealer
for the expected frame duration:

\begin{lstlisting}[language=Rust]
if data.len() <= 1 {
    return Ok(self.decode_lost(self.last_frame_size / self.ds));
}
\end{lstlisting}

\subsection{In-Band Forward Error Correction}\label{subsec:fec}

Concealment hides a loss but cannot recover the lost content. Opus can do better
for speech through \termdef{in-band forward error correction} (FEC): a SILK or
hybrid packet may carry a second, low-bitrate copy of the \emph{previous} frame,
called \termdef{LBRR} (low-bitrate redundancy). If a packet is lost but the next
one arrives, the decoder extracts the LBRR copy of the lost frame from the new
packet and decodes it instead of concealing, recovering a degraded but real
version of the audio. FEC is engaged through a separate decode entry point that,
on loss, prefers the LBRR data when present and conceals only when it is absent.
The cost is a modest bitrate increase on the carrying packets, which the encoder
budgets only when packet loss is expected.

\subsection{Discontinuous Transmission}\label{subsec:dtx}

During silence there is nothing to code, and continuing to send full frames
wastes bandwidth. \termdef{Discontinuous transmission} (DTX) lets the encoder
stop sending audio frames during inactivity, emitting only occasional one-byte
TOC packets that the decoder conceals as comfort noise. The decision
is driven by
an activity detector that tracks the background noise floor, so DTX engages
during genuine pauses even when there is audible ambient noise, not only during
digital silence:

\begin{lstlisting}[language=Rust]
fn frame_active(&mut self, pcm: &[f32]) -> bool {
    let ms = pcm.iter().map(|&v| v * v).sum::<f32>() / pcm.len() as f32;
    // Track the minimum energy with a slow upward leak (the noise floor).
    self.dtx_noise_floor = if ms < self.dtx_noise_floor { ms }
                           else { self.dtx_noise_floor * 1.0003 };
    ms > (self.dtx_noise_floor * ACTIVE_RATIO).max(ABS_FLOOR)
}
\end{lstlisting}

After a configured period of inactivity, the encoder drops to TOC-only packets,
cutting the silence bitrate to a fraction of a kilobit per second, and resumes
full coding the moment activity returns.

\subsection{The Encoder's Mode Selection}\label{subsec:encode_auto}

The encoder's top level chooses a mode per frame from the frame size and target
bitrate. Short frames (2.5 and 5~ms) can only be CELT; long frames (40 and
60~ms) can only be SILK; the 10 and 20~ms sizes admit all three, and the choice
falls to SILK for low-bitrate speech bandwidths, hybrid for the
higher bandwidths
at moderate rates, and CELT otherwise. Ahead of every mode, the input passes
through a 3~Hz one-pole high-pass that removes DC and sub-audible rumble before
coding, with its state carried across frames so the filter does not restart each
frame.

\section{Ogg Encapsulation}\label{sec:ogg}

The packet layer (\cref{sec:packet}) deliberately stores no length, assuming a
container provides framing. For stored files, that container is Ogg
(RFC 3533~\cite{rfc3533}), with an Opus-specific mapping (RFC
7845~\cite{rfc7845})
that defines the headers and timing. This section covers how Opus packets are
written to and read from an \texttt{.opus} file: the page structure, the
checksum, packet reassembly, the two mandatory headers, and the granule timing
that synchronizes audio to a clock.

\subsection{Pages and Lacing}\label{subsec:ogg_pages}

Ogg's transport unit is the \termdef{page}, not the packet. A page is a
self-contained, checksummed run of bytes that begins with the four-byte capture
pattern \code{OggS} and carries a header (flags, a granule position,
  the stream's
serial number, a page sequence number, and a CRC) followed by a \termdef{segment
table} and the segment data.

\begin{concept}{Lacing: Packets into Segments}
  A page does not store packet lengths directly. Instead its body is
  divided into
  \termdef{segments} of up to 255 bytes, and the segment table lists each
  segment's length as one byte. A packet is encoded as a sequence of 255-byte
  segments followed by one final segment of length less than 255, so a segment
  value below 255 marks the end of a packet and a value of exactly 255 means the
  packet continues into the next segment. A packet whose length is an exact
  multiple of 255 therefore ends with a zero-length segment, which is how the
  reader distinguishes ``the packet happened to fill the last segment'' from
  ``the packet continues.'' Because a page holds at most 255 segments, a packet
  too long for one page spills into the next, with a continuation
  flag marking the
  follow-on page.
\end{concept}

\subsection{The Ogg CRC}\label{subsec:ogg_crc}

Every page carries a 32-bit CRC over its entire contents (with the CRC field
itself zeroed during computation), so a reader can detect corruption and a
mistaken sync to a false \code{OggS} pattern. The polynomial is the standard
$0\texttt{x}04\text{C}11\text{DB}7$, but, unusually, Ogg uses the
\emph{non-reflected} variant with zero initial value and no final inversion,
which differs from the more common reflected CRC-32 of zlib or Ethernet.

\begin{lstlisting}[language=Rust]
const POLY: u32 = 0x04C1_1DB7;        // non-reflected, zero init, no final XOR

pub(crate) const fn update(mut crc: u32, data: &[u8]) -> u32 {
    let mut i = 0;
    while i < data.len() {
        crc = (crc << 8) ^ TABLE[((crc >> 24) as u8 ^ data[i]) as usize];
        i += 1;
    }
    crc
}
\end{lstlisting}

When a page fails its CRC, the reader does not give up: it scans forward for the
next \code{OggS} capture pattern and resynchronizes, so a single corrupt page
costs only its own audio, not the rest of the stream.

\subsection{Packet Reassembly}\label{subsec:ogg_reassembly}

Reading Opus packets from an Ogg file is the inverse of lacing: walk
the pages of
the chosen logical stream (matched by serial number), concatenate segments until
a terminator (a segment under 255 bytes) completes a packet, and carry a partial
packet across a page boundary when the page ends mid-packet. The reader tracks
the page sequence number to detect gaps, and discards a buffered partial packet
if continuity breaks, so a lost page cannot fuse two unrelated packets:

\begin{lstlisting}[language=Rust]
// A page ends mid-packet (last lacing == 255) -> buffer the partial; otherwise
// the accumulated bytes complete a packet. A sequence gap drops the partial.
if !consecutive || page.continued != self.have_partial {
    self.partial.clear();
    self.have_partial = false;
}
\end{lstlisting}

A defensive cap bounds the size of a reassembled packet, so hostile input cannot
drive unbounded allocation; a packet that exceeds it is dropped rather than
buffered.

\subsection{OpusHead and OpusTags}\label{subsec:ogg_headers}

An Ogg Opus stream begins with two mandatory header packets, each alone on its
own page. The first, \termdef{OpusHead}, is the identification
header: it carries
the version, the channel count, the \termdef{pre-skip} (\cref{subsec:granule}),
the original input sample rate (informational only, since decoding is always at
48~kHz internally), an output gain, and the channel mapping family that
determines whether the stream is mono/stereo or multichannel surround
(\cref{sec:multistream}). The second, \termdef{OpusTags}, is the comment header:
a vendor string and a list of \code{TAG=value} metadata entries, in the same
format as Vorbis comments.

\begin{lstlisting}[language=Rust]
pub struct OpusHead {
    pub version: u8,                    // 1
    pub channel_count: u8,
    pub pre_skip: u16,                  // 48 kHz samples to discard at startup
    pub input_sample_rate: u32,         // informational
    pub output_gain_q8: i16,            // Q7.8 dB, applied by the player
    pub channel_mapping: ChannelMapping,
}
\end{lstlisting}

\subsection{Granule Positions and Timing}\label{subsec:granule}

Ogg pages carry a \termdef{granule position}, a per-codec timestamp. For Opus it
counts samples at 48~kHz, the same reference clock used everywhere in the codec
(\cref{subsec:internal_rate}). The granule position of a page is the
total number
of decoded samples available after decoding every packet that
\emph{completes} on
that page, which lets a player seek and compute durations precisely.

\begin{concept}{Pre-Skip and End-Trim}
  Two timing offsets account for the codec's own latency. The encoder's filters
  (the SILK resampler, the CELT overlap) introduce a fixed delay, so the first
  decoded samples are startup transient rather than real audio. \termdef{Pre-
  skip}, stored in OpusHead, tells the player how many leading 48~kHz samples to
  discard so the output aligns with the original input. At the other end, the
  final page's granule position can fall short of the samples the last packet
  actually produces, signaling \termdef{end-trim}: the player discards the
  trailing samples beyond the granule position. Together these let an Opus file
  represent an audio signal of \emph{any} sample count exactly, even though the
  codec only produces audio in fixed-size frames. The reader resolves per-packet
  positions by walking backward from a page's granule position, subtracting each
  packet's duration (read from its TOC), so every packet gets a precise
  timestamp.
\end{concept}

The implementation provides both directions: a reader that parses headers,
reassembles packets, and resolves their granule positions, and a writer that
emits the two header pages, then audio pages with correctly computed granule
positions and a final end-of-stream page. The round trip is interoperable with
the reference tools: files written by the implementation decode in
ffmpeg/libopus at the correct duration, and files those tools produce decode
here.

\section{Multistream and Surround}\label{sec:multistream}

Opus codes mono or stereo. Surround content (5.1, 7.1, ambisonics) is
handled not
by extending the codec but by running several Opus streams in parallel and
multiplexing them into one packet, the \termdef{multistream} format
(RFC 7845~\cite{rfc7845} \S5.1.1). This keeps the core codec simple: it never
needs to know about more than two channels.

\subsection{Streams and Coupling}\label{subsec:ms_streams}

A multistream packet contains $N$ Opus streams. The first $M$ of them are
\termdef{coupled} (stereo, carrying two channels each) and the remaining $N-M$
are \termdef{uncoupled} (mono), so the streams together carry $2M + (N-M) = M+N$
decoded channels. A channel mapping table then routes those decoded channels to
the output layout.

\begin{concept}{Coupled and Uncoupled Streams}
  Pairing channels into a coupled (stereo) stream lets Opus exploit the
  correlation between them, exactly the mid/side coding of
  \cref{subsec:silk_stereo,subsec:celt_bands}, which is more efficient than
  coding two channels independently. Channels that are largely independent (a
  surround channel and the front left, say) are better coded as separate mono
  streams. An encoder partitions the output channels into coupled pairs and mono
  singletons, records that partition as the stream count $N$ and coupled count
  $M$, and stores a \termdef{channel mapping}: one entry per output channel
  giving which decoded channel feeds it, or a sentinel meaning silence. A
  5.1 layout, for instance, might use three coupled streams and no mono streams,
  with the mapping assigning the six decoded channels to front, center/LFE, and
  surround positions.
\end{concept}

\subsection{Demultiplexing}\label{subsec:ms_demux}

The streams are concatenated in one packet, which raises the framing problem the
core format avoided: the decoder must find where each stream's bytes end. Here
the \termdef{self-delimiting} framing from \cref{subsec:self_delimit} earns its
place. Every stream except the last is written self-delimited,
announcing its own
length, so the demultiplexer can parse them in order; the last stream is parsed
normally and consumes the remainder. Each stream is then decoded by its own Opus
decoder, and the results are scattered into the output according to the mapping:

\begin{lstlisting}[language=Rust]
for s in 0..streams {
    let packet = if s != streams - 1 {
        let (packet, used) = Packet::parse_self_delimited(rest)?;   // knows its length
        rest = &rest[used..];
        packet
    } else {
        Packet::parse(rest)?                                        // takes the rest
    };
    stream_pcm.push(self.decoders[s].decode_parsed(&packet));
}
// Route decoded channels to output channels (255 == silence).
for (ch, &m) in self.mapping.iter().enumerate() {
    if m == 255 { continue; }
    /* copy decoded channel m into output channel ch */
}
\end{lstlisting}

All streams in a packet must decode to the same number of samples, which the
demultiplexer enforces; a mismatch is rejected as malformed. The channel mapping
family in OpusHead (\cref{subsec:ogg_headers}) selects the convention: family 0
is plain mono/stereo, and the higher families define the surround channel orders
and the stream/coupling counts.

\section{Computational Optimizations}\label{sec:optimizations}

The preceding sections describe what Opus computes. This section describes how
the implementation makes those computations fast enough for real-time use,
decoding at hundreds of times realtime on one core, and encoding faster than the
reference at its default complexity. The optimizations fall into two groups:
those that are free because they change only encoder \emph{decisions}, and those
that preserve the bitstream exactly. Understanding which is which requires
returning to the asymmetry introduced in \cref{subsec:bits_overview}.

\subsection{The Licence to Approximate}\label{subsec:licence}

\begin{concept}{Why the Encoder May Approximate but the Decoder May Not}
  The Opus bitstream is defined by the sequence of range-coder symbols
  (\cref{sec:rangecoder}), not by the floating-point values an
  encoder computed to
  choose them. Any \emph{valid} choice the decoder can follow produces a legal
  stream. This splits the codec's arithmetic into two regimes with different
  rules. The decoder, and every encoder computation that feeds a coded symbol
  whose exact value the decoder recomputes (the bit allocation, the energy
  prediction, the NSQ reconstruction), must be \emph{bit-exact}: a single
  diverging bit desynchronizes the range coder. But an encoder computation that
  only \emph{informs a decision}, such as which pitch lag looks best or which
  pulse vector maximizes correlation, may use any approximation that lands on a
  valid answer, because the decoder never recomputes it; it just follows the
  symbols that result. Every SIMD kernel below lives in the second regime, which
  is why fast, approximate arithmetic is provably harmless there.
\end{concept}

This is also why the implementation can keep \code{unsafe} confined to a handful
of SIMD kernels. The crate denies \code{unsafe} by default; each kernel opts in
explicitly with a documented safety justification, and each sits on
the encoder's
decision path where the range coder, not the kernel's exact output, defines the
bitstream. Every round-trip and conformance test passes whether the SIMD path or
the scalar fallback runs.

\subsection{SIMD Dot Products}\label{subsec:simd_dot}

The dominant cost of encoding is correlation: the pitch analyses in
both SILK and
CELT, the LPC whitening filters, and the NSQ prediction all reduce to dot
products over long buffers. The implementation provides vectorized dot products
with runtime CPU-feature detection, dispatching to an AVX2+FMA kernel when
available, an SSE2 kernel as the x86-64 baseline, and a scalar fallback
elsewhere:

\begin{lstlisting}[language=Rust]
pub(crate) fn dot(x: &[f32], y: &[f32]) -> f32 {
    #[cfg(target_arch = "x86_64")]
    // SAFETY: both slices expose >= x.len() contiguous f32s; the kernels read
    // only those lanes (vector loads masked to width, with a scalar tail).
    unsafe { if has_avx2() { dot_avx2(x, y) } else { dot_sse2(x, y) } }
    #[cfg(not(target_arch = "x86_64"))]
    dot_scalar(x, y)
}
\end{lstlisting}

The feature detection is cached in an atomic so the check runs once,
not per call.
One subtlety bridges the two regimes of \cref{subsec:licence}: SILK's pitch
decisions are pinned against the reference, so the helper that feeds them
accumulates in an \code{f64} even though it operates on \code{f32} data, keeping
the SIMD result within the reference's tolerance. The CELT pitch
correlations and
the encoder's other analyses feed pure decisions and tolerate FMA's fused
rounding freely.

\subsection{The SIMD PVQ Search}\label{subsec:simd_pvq}

The pulse search of \cref{eq:pvq_search} is the hottest loop in CELT
encoding: for
each of $K$ pulses, it scans all $N$ positions for the one maximizing a
correlation ratio. The implementation vectorizes this with SSE2 and
AVX2 kernels,
and it is the clearest illustration of the licence to approximate.

\begin{concept}{Two Tricks: Reciprocal Square Roots and Sentinel Lanes}
  The correlation ratio of \cref{eq:pvq_search} needs a division by
  $\norm{\mathbf{y}}$, computed across lanes with the hardware
  \termdef{reciprocal square root} instruction, a fast 12-bit approximation
  rather than an exact divide. This can occasionally pick a different position
  than the scalar search would, but the result is always a \emph{valid} pulse
  vector, so it round-trips exactly through the decoder; only the
  encoder's choice
  shifts, never the bitstream's legality. The second trick handles the vector
  tail: the work buffers are padded to a whole number of SIMD lanes, and the
  padding lanes are filled with sentinel values (a very negative numerator and a
  very large denominator) that can never win the argmax. This lets every
  iteration process a full vector with no remainder branch and no bounds check,
  while guaranteeing a padding lane is never selected.
\end{concept}

The work buffers are thread-local and reused across bands and frames, so the
search allocates nothing in steady state. This pattern, reusing scratch buffers
rather than allocating per call, recurs across the hot paths: the PVQ
buffers, the
MDCT scratch, and the FFT plans are all cached per thread.

\subsection{The MDCT Pre-Rotation}\label{subsec:simd_mdct}

The MDCT's pre-rotation (\cref{subsec:mdct}) folds the real input into complex
values by a per-element complex multiply against precomputed twiddle factors.
Storing the twiddles in \termdef{folded} interleaved form (real and imaginary
parts adjacent) lets an SSE2 kernel deinterleave four complex
elements, multiply,
and reinterleave in a few instructions, replacing a scalar
complex-multiply loop.
Because the MDCT runs on the synthesis path as well, this kernel's
correctness is
guarded by the same TDAC perfect-reconstruction test that validates
the transform
itself (\cref{subsec:mdct}).

\subsection{The FFT Backend Seam}\label{subsec:fft_seam}

The largest single decode speedup comes not from a kernel but from a seam. The
MDCT's inner FFT (\cref{subsec:mdct}) is isolated behind a function
boundary with
two implementations selected at build time.

\begin{concept}{Two FFT Backends Behind One Interface}
  The dependency-free build evaluates the FFT directly, an $O(N^2)$ computation
  that keeps the crate free of dependencies for embedded and minimal-footprint
  targets. The default build instead routes the FFT through an external
  planned-FFT backend, an $O(N\log N)$ implementation, yielding
  roughly a tenfold
  decode speedup on the conformance vectors. The seam matters because the two
  backends produce subtly different floating-point results, and the codec's
  conformance criterion is PCM \emph{quality}, not bit-exact PCM
  (\cref{sec:numerical}): both backends pass, differing only in
  rounding far below
  the audibility threshold. Isolating the FFT this way lets the
  implementation be
  fast by default and dependency-free by choice, without forking the codec.
\end{concept}

\subsection{Table-Driven Enumeration and Buffer Reuse}\label{subsec:tables}

Two further optimizations are algorithmic rather than architectural. The CWRS
enumeration (\cref{subsec:cwrs}) replaces the $O(NK)$ recurrence with $O(N)$
lookups into a table computed once at compile time (a \code{const fn} baked
into the binary, not a runtime lazy-init), which is significant because band
decoding calls it on every band of every frame. And throughout the
hot paths, per-call heap allocation is eliminated by reusing
thread-local scratch
buffers: the band-encode work vectors, the MDCT complex scratch, and the cached
FFT plans all persist across calls, so steady-state coding does no allocation at
all. Combined with the \code{no\_std}-compatible core, this keeps the codec's
memory behavior predictable enough for real-time and embedded use.

\section{Numerical Considerations and Conformance}\label{sec:numerical}

Opus is specified with unusual precision about arithmetic, because much of it
must be reproduced exactly across implementations while the rest is deliberately
left free. This final section explains where bit-exactness is required, how
fixed-point and floating-point arithmetic divide the codec, and how an
implementation proves itself correct against the official conformance suite.

\subsection{Fixed-Point and Floating-Point}\label{subsec:fixed_float}

\begin{concept}{Q-Format Fixed-Point Arithmetic}
  A \termdef{\Q{n}} value stores a real number $x$ as the integer
  $\round{x \cdot 2^n}$, so the fraction lives in the low $n$ bits. Addition is
  ordinary integer addition; multiplying a \Q{a} value by a \Q{b} value gives a
  \Q{(a+b)} product that is shifted right to renormalize. Fixed-point arithmetic
  is exactly reproducible on every machine, with no dependence on a
  floating-point
  unit's rounding modes, which is why the normative parts of the
  codec use it. The
  cost is constant vigilance about scaling and overflow: the implementation
  matches the reference's wrapping integer semantics precisely, using explicit
  wrapping operations where the reference relies on machine wraparound, so that
  intermediate sums agree to the bit.
\end{concept}

The two coders make different choices. SILK is specified in fixed point, and the
implementation follows suit: its NLSF conversion, gain dequantization, LTP
filters, and NSQ run in \Q{}-format integers pinned bit-exactly against the
compiled reference, so SILK reconstructs sample-identical PCM. CELT's reference
exists in both fixed-point and floating-point forms, and RFC 6716 explicitly
permits a floating-point decoder; the implementation uses float for CELT, which
makes its PCM output differ from a fixed-point decoder by rounding far below
audibility, while its range-coder arithmetic remains exact.

\subsection{What Must Be Exact}\label{subsec:what_exact}

The dividing line follows the range coder. Anything that produces or consumes a
coded symbol must be bit-exact, because the symbols are the bitstream and the
range coder desynchronizes on any divergence. Anything downstream of the last
symbol, the final PCM samples after the inverse transform and resampling, need
only be perceptually correct.

\begin{importantnote}
  The single invariant that captures ``bit-exact'' is the final-range oracle of
  \cref{subsec:oracle}: after decoding a packet, the range coder's \code{rng}
  must equal the reference encoder's, exactly. This holds even for the
  floating-point CELT decoder, because the float arithmetic only affects the PCM
  \emph{after} the last symbol is decoded; every symbol along the way is decoded
  with identical integer range-coder arithmetic. An implementation can therefore
  use float for speed in the transform and still prove, packet by
  packet, that it
  decoded exactly the same bitstream the reference encoded.
\end{importantnote}

\subsection{The Conformance Criterion}\label{subsec:conformance}

RFC 6716 ships a suite of test vectors and a comparison tool,
\texttt{opus\_compare}, that scores a decoder's PCM output against the reference
on a perceptually weighted metric. An implementation conforms if it passes this
quality bar on every vector; bit-exact PCM is \emph{not} required, precisely
because the standard permits floating-point decoders whose output differs in
rounding.

The implementation discussed here passes the official criterion: all twelve
vectors score in the 99.2--100\% quality range, well above the passing bar, with
per-packet final ranges bit-exact across the whole suite. The combination is
strong evidence of correctness: the final-range equality proves the entropy
decode is exact symbol-for-symbol, and the quality score proves the signal
reconstruction is perceptually faithful. For the CELT-only vectors, the
synthesized PCM additionally scores 83--104~dB SNR against the reference float
decode, far beyond what the quality metric demands.

\begin{concept}{Two Checks, Two Purposes}
  The two conformance measures answer different questions. The
  \textbf{final-range oracle} answers ``did I decode the same bitstream?'', a
  binary, unforgiving check on the entropy coder and everything that feeds it; a
  single wrong symbol fails it. The \textbf{\texttt{opus\_compare} quality
  score} answers ``does the audio sound right?'', a graded check on the signal
  path that tolerates the rounding differences a floating-point decoder
  legitimately introduces. Passing both means an implementation is exact where
  the standard demands exactness and faithful where it demands
  fidelity, which is
  the strongest practical guarantee a decoder can offer.
\end{concept}

% ===END-OF-MANUAL===
